% Transcription of Cayley's Collected Mathematical Papers, Volume 4
% PDF pages 326--375 (printed pages 304--353)
% Papers: 267 (conclusion), 268, 269, 270, 271

\documentclass[11pt]{article}

% ============================================================
% PREAMBLE
% ============================================================
\usepackage[utf8]{inputenc}
\usepackage[T1]{fontenc}
\usepackage[english,french]{babel}
\usepackage{amsmath,amssymb}
\usepackage{geometry}
\usepackage{fancyhdr}
\usepackage{hyperref}
\usepackage{array}
\usepackage{longtable}
\usepackage[protrusion=true,expansion=false]{microtype}

\geometry{a4paper,margin=1in}

\pagestyle{fancy}
\fancyhf{}
\fancyhead[LE,RO]{\thepage}
\renewcommand{\headrulewidth}{0pt}

% Fallback providecommands
\providecommand{\mathbbm}[1]{\mathbb{#1}}
\providecommand{\sgn}{\operatorname{sgn}}

\setlength{\abovedisplayskip}{4pt plus 1pt minus 1pt}
\setlength{\belowdisplayskip}{4pt plus 1pt minus 1pt}

% Cayley delimiter convention shorthand
\newcommand{\Cdel}{\mid}

\hypersetup{
  pdftitle={Cayley Collected Papers Vol IV, pp.~304--353},
  pdfauthor={Arthur Cayley},
  hidelinks
}

\selectlanguage{english}

\usepackage{graphicx}
\emergencystretch=3em
\tolerance=600
\begin{document}

\thispagestyle{empty}
\begin{center}
\vspace*{2cm}
{\Large\scshape The Collected Mathematical Papers}\\[0.4em]
{\large of}\\[0.4em]
{\Large\scshape Arthur Cayley}\\[1.5em]
{\large Volume IV --- pages 304--353}\\[0.6em]
{\normalsize Paper 267 (conclusion), Papers 268, 269, 270, 271}\\[2em]
\vfill
\end{center}
\newpage

\setcounter{page}{304}

% ============================================================
% PAPER 267 (conclusion)
% ON THE PORISM OF THE IN-- AND CIRCUMSCRIBED POLYGON
% ============================================================

\noindent\textbf{304}\hfill\textsc{on the porism of the in--and--circumscribed polygon.}\hfill\textbf{[267}

\medskip

\noindent and comparing this with
\[
   \sqrt{1 + 4 b \xi + 4(c + b^2)\xi^2 + 4(d + 2 b c)\xi^3},
\]
we have
\[
\begin{aligned}
   b &= i + 2,\\
   c + b^2 &= i^2 + 6 i + \nu + 5,\\
   d + 2 b c &= 4 i\,(i^2 + 2 i + \nu + 1),
\end{aligned}
\]
and thence
\[
\begin{aligned}
   b &= i + 2,\\
   c &= 2 i + \nu + 1,\\
   d &= -\,2 i \nu - 2 i;
\end{aligned}
\]
and by means of these values, or by effecting the development in a different manner, we find
\[
B = \left\{\begin{aligned} &i\\ &{}+ 2,\end{aligned}\right.
\]
\[
C = \left\{\begin{aligned} &i\cdot 2\\ &{}+\nu + 1,\end{aligned}\right.
\]
\[
D = 2\bigl\{\,i\,(-\nu - 1)\bigr\},
\]
\[
E = \left\{\begin{aligned} &i^2\cdot 4\nu\\ &{}+ i\cdot 8\,(\nu + 1)\\ &{}- (\nu + 1)^2,\end{aligned}\right.
\]
\[
F = 2\left\{\begin{aligned} &i^2\cdot -4\nu\\ &{}+ i\,(\nu+1)(3\nu - 5)\\ &{}+ 2\,(\nu + 1)^2,\end{aligned}\right.
\]
\[
G = 2\left\{\begin{aligned} &i^3\cdot 8\nu\\ &{}+ i^2\cdot 24\nu\\ &{}+ i\cdot 4\nu\,(-3\nu + 5)\\ &{}+ (\nu+1)(\nu - 7),\end{aligned}\right.
\]
\[
H = 4\left\{\begin{aligned} &i^3\cdot -8\nu\\ &{}+ i^2\cdot -32\nu\\ &{}+ i^2\cdot 4\nu\,(3\nu - 11)\\ &{}+ i^2\cdot 16\,(3\nu - 1)\\ &{}+ i\cdot (\nu + 1)(\nu - 7)(-5\nu + 3)\\ &{}- 4\,(\nu + 1)^2\,(\nu - 3).\end{aligned}\right.
\]

\newpage
\setcounter{page}{305}

\noindent\textbf{267]}\hfill\textsc{on the porism of the in--and--circumscribed polygon.}\hfill\textbf{305}

\medskip

\noindent These values give
\[
\begin{aligned}
   C &= \left\{\begin{aligned}&2 i\\ &{}+\nu + 1,\end{aligned}\right.\hspace{2em}= [4],\\[4pt]
   &\quad\hspace{2em}= [5],\\[2pt]
   D &= 2\bigl\{\,i\,(-\nu - 1)\bigr\},\\[4pt]
   CE - D^2 &= \left\{\begin{aligned}&i^2\cdot 8\nu\\ &{}+ i\cdot 4(\nu+1)\\ &{}- 2\,(\nu + 1)^2,\end{aligned}\right.\hspace{2em}= [6],
\end{aligned}
\]
\[
\begin{vmatrix} C, & D, & E\\ D, & E, & F\\ E, & F, & G\end{vmatrix} = \left\{\begin{aligned}&i^3\cdot 64\nu^2\\ &{}+ i^2\cdot -32\nu\,(\nu+1)(\nu-1)\\ &{}+ i\cdot -16\nu\,(\nu+1)^2\\ &{}+ i\cdot -8\,(3\nu - 1)(\nu+1)^2\\ &{}+ i\cdot -4\,(\nu+1)^3\end{aligned}\right.= [4]\,[8].
\]
\[
\begin{vmatrix} D, & E\\ E, & F\\ F, & G\end{vmatrix} = i\,(-\nu - 1)\left\{\begin{aligned}&16 i^2 \nu\\ &{}+ (\nu+1)^2\end{aligned}\right.
\]

\medskip

\noindent It will be remarked that $[4] = i(-\nu - 1)$ breaks up into the factors $i$ and $\nu + 1$; and so $[6] = (2 i - \nu - 1)\{4\nu i^2 + (\nu+1)^2\}$ breaks up into the factors $2 i - \nu - 1$ and $4\nu i^2 + (\nu+1)^2$.

\medskip

\noindent It may be added that the developed expression of $[6]$ is
\[
\left\{\begin{aligned}&i^2\cdot 8\nu\\ &{}+ i\cdot 4(\nu+1)\\ &{}- 2\,(\nu+1)^2,\end{aligned}\right.
\]
so that the difference between this and $[5]$ is $i^2\cdot 4\,(\nu+1)^2$, which is $= d^2$; this agrees with the former result.

\medskip

\noindent M.~Mention has also given, but not in a developed form, the formul\ae{} for the enneagon and the endecagon, and the following formula for the decagon, viz.
\[
   [16 i^2 \nu + (\nu+1)^2]^2 + 16 i^2 \nu\,(\nu+1)^2\cdot[2 i^2(1 - \nu) - (\nu+1)^2]^2 = 0.
\]

\hfill C.~IV.\hfill 39

\newpage
\setcounter{page}{306}

\noindent\textbf{306}\hfill\textsc{on the porism of the in--and--circumscribed polygon.}\hfill\textbf{[267}

\medskip

\begin{center}\textsc{IV.}\\[0.3em]\textit{Considerations as to the form of relation, in the case of two circles, for Polygons of an odd and even number of sides respectively.}\end{center}

\medskip

\noindent The relation between the two conics, or condition for the existence of the polygon, is the same whatever point of the circumscribed conic is taken as an angle of the polygon. Take for an angle, a point of intersection of the two conics. Consider first the case of the triangle; if a point of intersection is taken as a vertex $A$ of the triangle, then the sides $AB$, $AC$ coincide in direction with the tangent at $A$ to the inscribed conic $U$, hence $B$ and $C$ coincide together at the point where this tangent meets the circumscribed conic $V$; $BC$ is therefore a tangent of $V$, but it is by hypothesis a tangent of $U$; hence for the triangle the relation between the inscribed conic $U$ and the circumscribed conic $V$ may be stated as follows: viz.\ a common tangent of $U$ and $V$ meets $V$ at a point of contact of a tangent of $U$.

\medskip

\noindent In like manner for the quadrangle, if $A$ be taken at a point of intersection, the sides $AB$ and $AD$ will coincide in direction with the tangent to $U$ at this point, consequently $B$ and $D$ must coincide at the point where this tangent meets $V$; hence also $CB$, $CD$, the two tangents to $U$ from the point $C$, must coincide, or $C$ must be a point of intersection of the conics $U$, $V$. In other words, the pole, with respect to the inscribed conic $U$, of a common chord $AC$ of the two conics must lie on the circumscribed conic $V$; this is therefore the condition for the quadrilateral.

\medskip

\noindent In the ordinary mode of drawing the figures, with two conics which do not intersect, the points and lines employed in the foregoing constructions are imaginary, but the conics may be so drawn that these points and lines are all real.

\medskip

\noindent In general, for a polygon of an odd number, $2n+1$, of sides, then starting from a point of intersection, the sides will coincide in pairs, viz.\ the first and last, second and last but one, and so on, the middle or $(n+1)$th side being a common tangent of the two conics. But for a figure of an even number, $2n$, of sides, then starting from a point of intersection of the two conics, the $n$th side will terminate at a second point of intersection, and then the same series of sides will be repeated in the reverse order, so that the sides will coincide in pairs, first and last, second and last but one, $n$th and $(n+1)$th. For a figure of an odd number of sides, the relation involves only a single point of intersection, but for a figure of an even number of sides, it involves two points of intersection.

\medskip

\noindent Now in the case of two circles, for a polygon of an odd number of sides, the same relation is obtained, whether we take as the point of intersection one of the actual points of intersection, or a circular point at infinity, and the relation $[2n+1] = 0$ does not break up into factors. And so for a polygon of an even number of sides, then taking for the two points of intersection, the two actual points of intersection, or the two circular points at infinity, we have one form of result; but taking for them an actual point of intersection and a circular point at infinity, we have a different form of result; and the equation $[2n] = 0$ does break up into factors.

\newpage
\setcounter{page}{307}

\noindent\textbf{267]}\hfill\textsc{on the porism of the in--and--circumscribed polygon.}\hfill\textbf{307}

\medskip

\noindent This is verified very simply in the case of the quadrangle. Taking for the two points of intersection the circular points at infinity, the line joining them is the line infinity, and its pole (with respect to the inscribed circle) is the centre of this circle; the relation therefore is that the centre of the inscribed circle lies on the circumscribed circle. But when this is the case, it is easy to see that the pole (with respect to the inscribed circle) of the radical axis, lies also on the circumscribed circle; this pole and the centre of the inscribed circle are in fact the extremities of a diameter of the circumscribed circle. The condition thus obtained is $R^2 - a^2 = 0$ (which is M.~Mention's condition $i = 0$). We have next to find the analytical relation when the pole (with respect to the inscribed circle), of the line joining one of the actual points of intersection with a circular point at infinity is a point on the circumscribed circle. This I effect as follows:---taking $z = 0$ as the equation of line infinity, if the origin be taken on the middle point of the radical axis, and if $x = 0$ be the radical axis, then the equations of the two circles may be taken to be
\[
\begin{aligned}
\text{Inscribed circle,} &\quad x^2 + y^2 - 2 l x z - l'\,z^2 = 0,\\
\text{Circumscribed circle,} &\quad x^2 + y^2 - 2 L x z - L'\,z^2 = 0,
\end{aligned}
\]
a circular point at infinity is
\[
   x : y : z = 1 : i : 0,\quad (i = \sqrt{-1}),
\]
an actual point of intersection is
\[
   x : y : z = 0 : \sqrt{l'} : 1.
\]
The line joining these is
\[
   i x - y + z\sqrt{l'} = 0,
\]
its pole, with respect to the inscribed circle, is
\[
   x : y : z = -\,l\sqrt{l'} : i\sqrt{l'} - i l : 1;
\]
and if this be a point of the circumscribed circle
\[
   -\,l'\,(l' - 2 i l\sqrt{l'}) + 2 L i\sqrt{l'} - L' = 0,
\]
that is
\[
   2\,(L - l)\,i\,\sqrt{l'} = l'^2 + L',
\]
or
\[
   (l'^2 + L')^2 + 4\,(L - l)^2\,l' = 0,
\]
which is the required relation: but to express it in terms of the ordinary data $R, r, a$, the equations of the circles, putting therein $z = 1$, become
\[
\begin{aligned}
   (x - l)^2 + y^2 &= l' + l^2,\\
   (x - L)^2 + y^2 &= L' + L^2,
\end{aligned}
\]
\hfill 39--2

\newpage
\setcounter{page}{308}

\noindent\textbf{308}\hfill\textsc{on the porism of the in--and--circumscribed polygon.}\hfill\textbf{[267}

\medskip

\noindent and therefore
\[
\begin{aligned}
   a &= L - l,\\
   R^2 &= L' + L^2;
\end{aligned}
\]
these equations give
\[
\begin{aligned}
   r^2 &= l' + l^2,\\
   a &= L - l;
\end{aligned}
\]
and thence
\[
   L = \frac{R^2 - r^2 + a^2}{2a},
\]
\[
   l' = R^2 - \left(\frac{R^2 - r^2 + a^2}{2a}\right)^2 = \frac{1}{4 a^2}\bigl(2 r^2 R^2 + 2 r^2 a^2 + 2 R^2 a^2 - r^4 - R^4 - a^4\bigr);
\]
if with M.~Mention we write
\[
   \frac{1}{4 a^2}\,\bigl(r^4 + R^4 + a^4 - 2 r^2 R^2 - 2 r^2 a^2 - 2 R^2 a^2\bigr) = -\,\nu'.
\]

\medskip

\noindent The equation
\[
   (l'^2 + L')^2 + 4\,(L - l)^2\,l' = 0
\]
thus becomes
\[
   r^4 + 4 a^2\cdot\frac{-\,r^4}{4 a^2}\,\nu' = 0,
\]
that is, it becomes $\nu' + 1 = 0$, which is the other factor of the complete condition $\nu'(\nu' + 1) = 0$.

\newpage
\setcounter{page}{309}

% ============================================================
% PAPER 268
% ON A NEW AUXILIARY EQUATION IN THE THEORY OF
% EQUATIONS OF THE FIFTH ORDER
% ============================================================

\begin{center}
\textbf{268.}\\[1em]
{\Large ON A NEW AUXILIARY EQUATION IN THE THEORY OF}\\[0.3em]
{\Large EQUATIONS OF THE FIFTH ORDER.}\\[1em]
\textit{[From the Philosophical Transactions of the Royal Society of London, vol.~\textsc{cli}.\ (for the year 1861), pp.~263--276.\ \ Received February 20,---Read March 7, 1861.]}
\end{center}

\medskip

\noindent\textsc{Considering} the equation of the fifth order, or quintic equation,
\[
   (*\Cdel v,\,1)^5 = (v - x_1)(v - x_2)(v - x_3)(v - x_4)(v - x_5) = 0,
\]
and putting as usual
\[
   f\omega = x_1 + \omega x_2 + \omega^2 x_3 + \omega^3 x_4 + \omega^4 x_5,
\]
where $\omega$ is an imaginary fifth root of unity, then, according to Lagrange's general theory for the solution of equations, $f\omega$ is the root of an equation of the order 24, called the Resolvent Equation, but the solution whereof depends ultimately on an equation of the sixth order, viz.
\[
   (f\omega)^5,\ (f\omega^2)^5,\ (f\omega^3)^5,\ (f\omega^4)^5
\]
are the roots of an equation of the fourth order, each coefficient whereof is determined by an equation of the sixth order; and moreover the other coefficients can be all of them rationally expressed in terms of any one coefficient assumed to be known; the solution thus depends on a single equation of the sixth order. In particular the last coefficient, or
\[
   (f\omega\cdot f\omega^2\cdot f\omega^3\cdot f\omega^4)^5,
\]
is determined by an equation of the sixth order; and not only so, but its fifth root, or
\[
   f\omega\cdot f\omega^2\cdot f\omega^3\cdot f\omega^4,
\]
(which is a rational function of the roots, and is the function called by Mr~Cockle the Resolvent Product), is also determined by an equation of the sixth order: this

\newpage
\setcounter{page}{310}

\noindent\textbf{310}\hfill\textsc{on a new auxiliary equation in}\hfill\textbf{[268}

\medskip

\noindent equation may be called the Resolvent-Product Equation. But the recent researches of Mr~Cockle and Mr~Harley\footnote{Cockle, ``Researches in the Higher Algebra,'' \textit{Manchester Memoirs}, t.~\textsc{xv}.\ pp.~131--142 (1858). \ Harley, ``On the Method of Symmetric Products, and its Application to the Finite Algebraic Solution of Equations,'' \textit{Manchester Memoirs}, t.~\textsc{xv}.\ pp.~172--219 (1859).\ \ Harley, ``On the Theory of Quintics,'' \textit{Quart.\ Math.\ Journ.}\ t.~\textsc{iii}.\ pp.~343--359 (1859).} show that the solution of an equation of the fifth order may be made to depend on an equation of the sixth order, originating indeed in, and closely connected with, the resolvent-product equation, but of a far more simple form; this is the auxiliary equation referred to in the title of the present memoir. The connexion of the two equations, and the considerations which led to the new one, will be pointed out in the sequel; but I will here state synthetically the construction of the auxiliary equation. Representing for shortness the roots $(x_1, x_2, x_3, x_4, x_5)$ of the given quintic equation by $1, 2, 3, 4, 5$, and putting moreover
\[
   \overline{12345} = 12 + 23 + 34 + 45 + 51,\ \&c.
\]
(where on the right-hand side $12, 23, \&c.$ stand for $x_1 x_2,\ x_2 x_3,\ \&c.$), then the auxiliary equation, say
\[
   (*\Cdel\phi,\,1)^6 = 0,
\]
has for its roots
\[
\begin{aligned}
   \phi_1 &= \overline{12345} - \overline{24135},& \phi_4 &= \overline{21435} - \overline{13245},\\
   \phi_2 &= \overline{13425} - \overline{32145},& \phi_5 &= \overline{31245} - \overline{14325},\\
   \phi_3 &= \overline{14235} - \overline{43125},& \phi_6 &= \overline{41325} - \overline{12435},
\end{aligned}
\]
and, it follows therefrom, is of the form
\[
   (1,\ 0,\ C,\ 0,\ E,\ F,\ G\Cdel\phi,\,1)^6
\]
where $C, E, G$ are rational and integral functions of the coefficients of the given equation, being in fact seminvariants, and $F$ is a mere numerical multiple of the square root of the discriminant.

\medskip

\noindent The roots of the given quintic equation are each of them rational functions of the roots of the auxiliary equation, so that the theory of the solution of an equation of the fifth order appears to be now carried to its extreme limit. We have in fact
\[
\begin{aligned}
   \phi_1\phi_2 + \phi_3\phi_4 + \phi_5\phi_6 &= (*\Cdel x_2,\,1)^4,\\
   \phi_1\phi_5 + \phi_2\phi_3 + \phi_4\phi_6 &= (*\Cdel x_3,\,1)^4,\\
   \phi_1\phi_3 + \phi_2\phi_6 + \phi_3\phi_5 &= (*\Cdel x_4,\,1)^4,\\
   \phi_1\phi_4 + \phi_2\phi_5 + \phi_3\phi_6 &= (*\Cdel x_5,\,1)^4,
\end{aligned}
\]
where $(*\Cdel x_1,\,1)^4$, \&c.\ are the values, corresponding to the roots $x_1$, \&c.\ of the given equation, of a given quartic function. And combining these equations respectively with the quintic equations satisfied by the roots $x_1$, \&c.\ respectively, it follows that, conversely, the roots $x_1, x_2$, \&c.\ are rational functions of the combinations $\phi_1\phi_2 + \phi_3\phi_4 + \phi_5\phi_6$, $\phi_1\phi_5 + \phi_2\phi_3 + \phi_4\phi_6$, \&c.\ respectively, of the roots of the auxiliary equation.

\newpage
\setcounter{page}{311}

\noindent\textbf{268]}\hfill\textsc{the theory of equations of the fifth order.}\hfill\textbf{311}

\medskip

\noindent It is proper to notice that, combining together in every possible manner the six roots of the auxiliary equation, there are in all fifteen combinations of the form $\phi_1\phi_2 + \phi_3\phi_4 + \phi_5\phi_6$. But the combinations occurring in the above-mentioned equations are a completely determinate set of five combinations: the equation of the order 15, whereon depend the combinations $\phi_1\phi_2 + \phi_3\phi_4 + \phi_5\phi_6$, is not rationally decomposable into three quintic equations, but only into a quintic equation having for its roots the above-mentioned five combinations, and into an equation of the tenth order, having for its roots the other ten combinations, and being an irreducible equation. Suppose that the auxiliary equation and its roots are known; the method of ascertaining what combinations of roots correspond to the roots of the quintic equation would be to find the rational quintic factor of the equation of the fifth order, and observe what combinations of the roots of the auxiliary equation are also roots of this quintic factor. The direct calculation of the auxiliary equation by the method of symmetric functions would, I imagine, be very laborious. But the coefficients are seminvariants, and the process explained in my memoir on the Equation of Differences, [262], was therefore applicable, and by means of it, the equation, it will be seen, is readily obtained. The auxiliary equation gives rise to a corresponding covariant equation, which is given at the conclusion of the memoir.

\medskip

\noindent 1. \ I will commence by referring to some of the results obtained by Mr~Cockle and Mr~Harley.

\medskip

\noindent In the paper ``Researches on the Higher Algebra,'' Mr~Cockle, dealing with the quintic equation
\[
   v^5 - 5 Q v + E = 0,
\]
obtains for the Resolvent Product $\theta\,(= f\omega\,f\omega^2\,f\omega^3\,f\omega^4)$ the equation
\[
   \theta^6 + 2 Q E\cdot 5^2\,\theta^4 + 2 Q^4\cdot 5^4\,\theta^3 + Q^2 E^2\cdot 5^{10}\,\theta^2 - (5^8 Q^5 - E^2)\,5^4\,\theta + 5^{15} Q^6 = 0;
\]
and he remarks that this equation may be written
\[
   (\theta^3 + 5^2 Q E\theta + 5^7 Q^4)^2 = 5^{10}\,(108 Q^5 - E^2)\,\theta,
\]
so that $\sqrt[3]{-\theta}$ is determined by an equation of the sixth order, involving the quadratic radical $\sqrt{E(E^2 - 108 Q^5)}$, which is in fact the square root of the discriminant of the quintic equation.

\medskip

\noindent 2. \ Mr~Harley, in his paper ``On the Symmetric Product \&c.,'' makes use of the functions
\[
   \tau = x_1 x_2 + x_2 x_3 + x_3 x_4 + x_4 x_5 + x_5 x_1\ (= \overline{12345}),
\]
\[
   \tau' = x_1 x_3 + x_3 x_5 + x_5 x_2 + x_2 x_4 + x_4 x_1\ (= \overline{24135}),
\]
and he obtains for the form $v^5 - 5 Q v + E = 0$, the relation $\theta = 5\tau\tau'$, which, since here $\tau + \tau' = 0$, gives $\theta = -\,5\tau^2$.

\newpage
\setcounter{page}{312}

\noindent\textbf{312}\hfill\textsc{on a new auxiliary equation in}\hfill\textbf{[268}

\medskip

\noindent Hence $\tau\,(=\sqrt{-\tfrac{1}{5}\theta})$ is the root of an equation of the sixth order involving the radical $\sqrt{E\,(E^2 - 108 Q^5)}$, and which is in fact ($t = \tau\sqrt{5} = \sqrt{-\theta}$), the equation
\[
   t^6 + 5 Q E\,t^2 + \sqrt{E\,(E^2 - 108 Q^5)}\,t - 5 Q^4 = 0,
\]
given in Mr~Harley's paper ``On the Theory of Quintics.''

\medskip

\noindent 3. \ And in the same paper there is given a system of equations
\[
   t_1 t_2 + t_3 t_4 + t_5 t_6 = x_1\,(3 Q - x_1^2),\ \&c.,
\]
connecting the five roots of the given quintic equation with the combinations
\[
   t_1 t_2 + t_3 t_4 + t_5 t_6,\ \&c.
\]
of the roots of the equation in $t$.

\medskip

\noindent 4. \ I quote also, with a slight change of notation, the following results from the paper ``On the Symmetric Product \&c.,'' viz.\ considering the quintic equation under the form
\[
   (a, b, c, d, e, f\Cdel v,\,1)^5 = 0,
\]
we have
\[
\begin{aligned}
   f\omega\cdot f\omega^4 &= \Sigma x^2 + \tau\,(\omega + \omega^4) + \tau'\,(\omega^2 + \omega^3),\\
   f\omega^2\cdot f\omega^3 &= \Sigma x^2 + \tau\,(\omega^2 + \omega^3) + \tau'\,(\omega + \omega^4),
\end{aligned}
\]
where
\[
   \Sigma x^2 = x_1^2 + x_2^2 + x_3^2 + x_4^2 + x_5^2 = \frac{1}{a^2}\,(b^2 - 2 a c),
\]
and thence, observing also that $\tau + \tau' = -\dfrac{c}{a}$,
\[
   a^4 e\ (= a^4 f\omega\,f\omega^2\,f\omega^3\,f\omega^4) = 5 a^2 b d - 5 a b^2 c + b^4 + 5 a^2 \tau\tau',
\]
or, as this equation may also be written,
\[
   4 a^4 \theta = (5 a c - 2 b^2)^2 - 5 a^2\,(\tau - \tau')^2;
\]
and hence the Resolvent Product $\theta\,(= f\omega\,f\omega^2\,f\omega^3\,f\omega^4)$ being determined by an equation of the sixth order, this is also the case with the function $(\tau - \tau')^2$.

\medskip

\noindent 5. \ But the twelve functions $\pm\,(\tau - \tau')$ can be divided into two sets of six functions each, so that each set is determined by an equation of the sixth order involving a single quadratic radical. This was in fact suggested to me by Mr~Harley's equation in $t$; for in the case considered $\tau + \tau'$ was $= 0$, or $2\tau = \tau - \tau'$, and the equation in $t$ was presumably the particular form of the equation for $\tfrac{1}{2}(\tau - \tau')$ in the general case. But it will presently appear in what manner the conclusion should have been arrived at \textit{\`a priori}.

\newpage
\setcounter{page}{313}

\noindent\textbf{268]}\hfill\textsc{the theory of equations of the fifth order.}\hfill\textbf{313}

\medskip

\noindent 6. \ The preceding remarks show the connexion between the function $\phi\,(= \tau - \tau')$ to which belongs the new auxiliary equation, and the Resolvent Product $\theta\,(= f\omega\,f\omega^2\,f\omega^3\,f\omega^4)$. The relation was given for the denumerate form of the quintic; but taking, instead, the standard form $(a, b, c, d, e, f\Cdel v,\,1)^5 = 0$, it becomes
\[
   4 a^4 \theta = 2500\,(a c - b^2)^2 - 5 a^2\phi^2.
\]

\medskip

\noindent 7. \ The foregoing equation shows that $\phi$ is a seminvariantive function of the roots. In fact
\[
   f\omega = x_1 - x_5 + \omega\,(x_2 - x_5) + \omega^2\,(x_3 - x_5) + \omega^3\,(x_4 - x_5),
\]
is seminvariantive, and $f\omega^2$, $f\omega^3$, $f\omega^4$, being in like manner seminvariantive, the product $\theta\,(= f\omega\,f\omega^2\,f\omega^3\,f\omega^4)$ is also seminvariantive; $a c - b^2$ and $a$ are seminvariants, and therefore $\phi$ is a seminvariantive function.

\medskip

\noindent 8. \ But it is easy to show this directly. For representing, as before, the roots by $1, 2, 3, 4, 5$, we have
\[
\begin{aligned}
&(1-5)(2-5) + (2-5)(3-5) + (3-5)(4-5) = 12 + 23 + 34 - 5(1 + 2\cdot 2 + 2\cdot 3 + 4) + 3\cdot 5^2,\\
&(2-5)(4-5) + (4-5)(1-5) + (1-5)(3-5) = 24 + 41 + 13 - 5(2 + 2\cdot 4 + 2\cdot 1 + 3) + 3\cdot 5^2;
\end{aligned}
\]
and the difference of the right-hand sides is
\[
   12 + 23 + 34 - 24 - 41 - 13 + 5\,(2 + 4 + 3 - 4 - 1 - 3) - 0,
\]
which is $= \overline{12345} - \overline{24135}$. So that $\phi_1$,
\[
   = (1-5)(2-5) + (2-5)(3-5) + (3-5)(4-5) - [(2-5)(4-5) + (4-5)(1-5) + (1-5)(3-5)],
\]
is a function of the differences of the roots, that is, it is a seminvariantive function.

\medskip

\noindent 9. \ To account for the division of the twelve values of $\pm\,(\tau - \tau')$ into two sets as above, and to explain the formation of a set, consider the symbols $1, 2, 3, 4, 5$ as belonging to five points. We may with these five points form in all $(\tfrac{1}{2}\cdot 1\cdot 2\cdot 3\cdot 4 =)\,12$ pentagons, and the symbol $12345$ of any pentagon may of course be read backwards or forwards from any point ($12345 = 23451 = \&c.\ = 15432 = \&c.$) without alteration of its meaning. Now attaching to each arrangement of the five numbers a sign, $+$ or $-$, according to the ordinary rule of signs, $12345$ being as usual positive, the arrangements $12345$, $23451$, \&c.\ ${}\ldots\,15432$, \&c., which belong to the same pentagon, have all of them the same sign; and we may consequently connect with each pentagon the sign $+$ or $-$; there are, in fact, six pentagons with the sign $+$ and six with the sign $-$; and to each positive pentagon there corresponds a negative pentagon, which is derived from it by stellation, viz.\ to the positive pentagon $12345$ there corresponds the negative one $24135$, and so for the other positive pentagons. The above-mentioned system of equations
\[
\begin{aligned}
   \phi_1 &= \overline{12345} - \overline{24135},& \phi_4 &= \overline{21435} - \overline{13245},\\
   \phi_2 &= \overline{13425} - \overline{32145},& \phi_5 &= \overline{31245} - \overline{14325},\\
   \phi_3 &= \overline{14235} - \overline{43125},& \phi_6 &= \overline{41325} - \overline{12435},
\end{aligned}
\]

\hfill C.~IV.\hfill 40

\newpage
\setcounter{page}{314}

\noindent\textbf{314}\hfill\textsc{on a new auxiliary equation in}\hfill\textbf{[268}

\medskip

\noindent in fact exhibits the six positive pentagons, each accompanied by its stellated negative pentagon, and the formation of the system of equations is thus completely explained; the order of arrangement of the pairs \textit{inter se} (or, what comes to the same thing, the order of arrangement of the suffixes of the $\phi$'s) is wholly immaterial.

\medskip

\noindent 10. \ The six pairs of pentagons, or, what is the same thing, the $\phi$'s, correspond to each other in pairs in a fivefold manner, \textit{quoad} the numbers $1, 2, 3, 4, 5$ respectively; thus, \textit{quoad} 5, the pairs are $\phi_1$ and $\phi_4$, $\phi_2$ and $\phi_5$, $\phi_3$ and $\phi_6$, or say $1$ and $4$, $2$ and $5$, $3$ and $6$. The relation is best seen by means of the positive pentagons; thus, \textit{quoad} 5, in the pentagons $12345$ and $21435$, the points adjacent to $5$ in the one of them are the points $2, 3$, and in the other of them the complementary points $1, 4$; and so in the other cases. The fivefold correspondence is shown by the symbolical equations
\[
\begin{aligned}
   1 &= 16,\ 24,\ 35,\\
   2 &= 12,\ 34,\ 56,\\
   3 &= 15,\ 23,\ 46,\\
   4 &= 13,\ 26,\ 35,\\
   5 &= 14,\ 25,\ 36,
\end{aligned}
\]
which, in fact, indicate the combinations of the $\phi$'s which correspond to the several roots of the quintic.

\medskip

\noindent 11. \ It is proper to notice that the right-hand sides of the last-mentioned equations contain all the duads formed with the six numbers $1, 2, 3, 4, 5, 6$, each duad once, and once only. There are in all six such synthemes of duads, viz.
\[
\begin{array}{lll}
12.34.56 & 12.35.46 & 12.36.45\\
13.25.46 & 13.24.56 & 13.25.46\\
14.26.35 & 14.25.36 & 14.23.56\\
15.24.36 & 15.26.34 & 15.26.34\\
16.23.45 & 16.23.45 & 16.24.35\\[6pt]
12.34.56 & 12.35.46 & 12.36.45\\
13.26.45 & 13.26.45 & 13.24.56\\
14.25.36 & 14.23.56 & 14.26.35\\
15.23.46 & 15.24.36 & 15.23.46\\
16.24.35 & 16.25.34 & 16.25.34
\end{array}
\]
which is in fact the theorem whereon depends the existence, for six letters, of a 6-valued function not symmetrical in respect of five letters. There is not any peculiarity in the syntheme of duads which above presented itself; the occurrence of this particular syntheme, instead of any other, arises merely from the arbitrary selection of the suffixes of the $\phi$'s.

\medskip

\noindent 12. \ It is hardly necessary to remark that if the pentagon $12345$ had been assumed negative instead of positive, the only difference would be that the $\phi$'s would have their signs reversed.

\newpage
\setcounter{page}{315}

\noindent\textbf{268]}\hfill\textsc{the theory of equations of the fifth order.}\hfill\textbf{315}

\medskip

\noindent 13. \ I proceed now to the calculation of the Auxiliary Equation. As the working is rather easier for that form, I shall in the first instance take for the given quintic the denumerate form
\[
   (a, b, c, d, e, f\Cdel v,\,1)^5 = 0.
\]
Representing, as before, the roots $x_1, x_2, x_3, x_4, x_5$ of this equation by $1, 2, 3, 4, 5$, and writing
\[
   \overline{12345} = 12 + 23 + 34 + 45 + 51,\ \&c.
\]
(where on the right-hand side $12, 23, \&c.$ stand for $x_1 x_2,\ x_2 x_3,\ \&c.$), we have to find the equation
\[
   (*\Cdel\phi,\,1)^6 = 0,
\]
the roots whereof are
\[
\begin{aligned}
   \phi_1 &= \overline{12345} - \overline{24135},& \phi_4 &= \overline{21435} - \overline{13245},\\
   \phi_2 &= \overline{13425} - \overline{32145},& \phi_5 &= \overline{31245} - \overline{14325},\\
   \phi_3 &= \overline{14235} - \overline{43125},& \phi_6 &= \overline{41325} - \overline{12435}.
\end{aligned}
\]
As already remarked, the coefficients are seminvariants, and if the equation is in the first instance calculated for the particular case $f = 0$, the terms in $f$ can be separately determined. But putting $f = 0$, one of the roots, say $5$, becomes $= 0$, and the remaining roots $1, 2, 3, 4$ are the roots of the quartic equation $(a, b, c, d, e\Cdel v,\,1)^4 = 0$.

\medskip

\noindent 14. \ Writing for shortness
\[
   \overline{1234} = 12 + 23 + 34,\ \&c.,
\]
and putting also
\[
\begin{aligned}
   A &= 12 + 34,\\
   B &= 13 + 42,\\
   C &= 14 + 23,
\end{aligned}
\]
then we have
\[
\begin{aligned}
   \phi_1 &= \overline{1234} - \overline{2413} = 12 + 23 + 34 - 24 - 41 - 13 = A - B + 23 - 14,\\
   \phi_2 &= \overline{1342} - \overline{3214} = 13 + 34 + 42 - 32 - 21 - 14 = B - C + 34 - 12,\\
   \phi_3 &= \overline{1423} - \overline{4312} = 14 + 42 + 23 - 43 - 31 - 12 = C - A + 42 - 13,\\
   \phi_4 &= \overline{2143} - \overline{1324} = 21 + 14 + 43 - 13 - 32 - 24 = A - B - 23 + 14,\\
   \phi_5 &= \overline{3124} - \overline{1423} = 31 + 12 + 24 - 14 - 42 - 23 = B - C - 34 + 12,\\
   \phi_6 &= \overline{4132} - \overline{1243} = 41 + 13 + 32 - 12 - 24 - 43 = C - A - 42 + 13.
\end{aligned}
\]

\medskip

\noindent 15. \ We have then
\[
   (\phi - \phi_1)(\phi - \phi_4) = (\phi - A + B)^2 - (14 - 23)^2,
\]
\hfill 40--2

\newpage
\setcounter{page}{316}

\noindent\textbf{316}\hfill\textsc{on a new auxiliary equation in}\hfill\textbf{[268}

\medskip

\noindent where $\overline{1234}$ denotes the product of the four roots; the functions $A, B, C$, and the product $\overline{1234}$, are each of the degree zero in the coefficients $(a, b, c, d, e)$; and if we put
\[
\begin{aligned}
   b &= -\,c',\\
   c &= -\,4 a e + b d,
\end{aligned}
\]
then we in fact have
\[
\begin{aligned}
   d &= 4 a c e - a d^2 - b^2 e,\\
   a\,\Sigma A &= -\,b,\\
   a^2\,\Sigma A B &= c,\\
   a^3\,A B C &= -\,d,\\
   a^4\cdot\overline{1234} &= e.
\end{aligned}
\]

\medskip

\noindent But on the understanding that $\phi$ is ultimately to be changed into $a\phi$, it is allowable, and it will be convenient to write
\[
\begin{aligned}
   \Sigma A &= -\,b,\\
   \Sigma A B &= c,\\
   A B C &= -\,d,\\
   \overline{1234} &= a e.
\end{aligned}
\]

\medskip

\noindent 16. \ I assume also
\[
\begin{aligned}
   B + C - A &= \alpha,\\
   C + A - B &= \beta,\\
   A + B - C &= \gamma;
\end{aligned}
\]
we have thus
\[
\begin{aligned}
   (\phi - \phi_1)(\phi - \phi_4) &= (\phi + \alpha)(\phi - \beta) + 4 a e,\\
   (\phi - \phi_2)(\phi - \phi_5) &= (\phi + \beta)(\phi - \gamma) + 4 a e,\\
   (\phi - \phi_3)(\phi - \phi_6) &= (\phi + \gamma)(\phi - \alpha) + 4 a e,
\end{aligned}
\]
so that the equation in $\phi$ is
\[
   [(\phi + \beta)(\phi - \gamma) + 4 a e]\,[(\phi + \gamma)(\phi - \alpha) + 4 a e]\,[(\phi + \alpha)(\phi - \beta) + 4 a e] = 0.
\]

\medskip

\noindent 17. \ To obtain the symmetrical functions of $\alpha, \beta, \gamma$ it is only necessary to remark that if in the identical equation
\[
   (1, b, c, d\Cdel\theta,\,1)^3 = (\theta - A)(\theta - B)(\theta - C),
\]
we put $\tfrac{1}{2}(x + A + B + C) = \tfrac{1}{2}(x - b)$, in the place of $\theta$, the equation becomes
\[
   (1, b, c, d\Cdel\tfrac{1}{2}(x - b),\,2)^3 = (x + \alpha)(x + \beta)(x + \gamma),
\]
so that we have
\[
\begin{aligned}
   \Sigma\alpha &= -\,b\quad =\quad c,\\
   \Sigma\alpha\beta &= -\,b^2 + 4 c = -\,16 a e + 4 b d - c^2,\\
   \alpha\beta\gamma &= b^3 - 4 b c + 8 d = 16 a c e - 4 a d^2 - 8 b^2 e + 4 b c d - c^3.
\end{aligned}
\]

\newpage
\setcounter{page}{317}

\noindent\textbf{268]}\hfill\textsc{the theory of equations of the fifth order.}\hfill\textbf{317}

\medskip

\noindent 18. \ The developed expression for the equation in $\phi$ is easily found to be
\[
\begin{aligned}
   &\phi^6\\
   &{} + \phi^4\cdot(-\,2\,\Sigma\alpha^2\quad + 12 a e)\\
   &{} + \phi^3\cdot(2\,\Sigma\alpha^2\beta^2 - 4 a e\,(2\,\Sigma\alpha^2 + \Sigma\alpha\beta) + 48 a^2 e^2)\\
   &{} + \phi\cdot(-\,4 a e\,(\alpha - \beta)(\beta - \gamma)(\gamma - \alpha))\\
   &{} + (-\,\alpha^2\beta^2\gamma^2 + 4 a e\,\alpha\beta\gamma\,\Sigma\alpha - 16 a^2 e^2\,\Sigma\alpha\beta + 64 a^3 e^3) = 0.
\end{aligned}
\]

\medskip

\noindent 19. \ In this equation the coefficient of $\phi$ is
\[
\begin{aligned}
   &-\,4 a e\,\Sigma(B - A)(C - B)(A - C)\\
   &{} = 32 a e\,(A - B)(B - C)(C - A);
\end{aligned}
\]
or, neglecting the multiplier $a$, it is
\[
   -\,32\cdot 1\cdot 2\cdot 3\cdot 4\,(1 - 2)(1 - 3)(1 - 4)(2 - 3)(2 - 4)(3 - 4),
\]
which is the value for $a = 0$, of
\[
   -\,32\,(1 - 2)(1 - 3)(1 - 4)(1 - 5)(2 - 3)(2 - 4)(2 - 5)(3 - 4)(3 - 5)(4 - 5),
\]
i.e.\ the coefficient in question is
\[
   -\,32\cdot 25\,\sqrt{5}\,\sqrt{a^2 f^2 + \&c.} = -\,800\,\sqrt{5}\,\sqrt{a^2 f^2 + \&c.},
\]
where $a^2 f^2 + \&c.$ denotes the discriminant of the denumerate form
\[
   (a, b, c, d, e, f\Cdel v,\,1)^5.
\]

\medskip

\noindent 20. \ The remaining coefficients are rational functions of $a, b, c, d, e$, which have to be completed by the introduction of the terms in $f$. We have
\[
\begin{aligned}
   \text{Coeff.\ }\phi^4 &= -\,(\Sigma\alpha)^2 + 2\,(-\,16 a e + 4 b d - c^2)\\
   &\qquad{} + 2\,\Sigma\alpha\beta + 12 a e.
\end{aligned}
\]
\[
\begin{aligned}
   \text{Coeff.\ }\phi^3 &= (-\,16 a e + 4 b d - c^2)^2\\
   &\qquad{} - 2\,(a^2\beta^2\gamma^2)\Sigma\alpha - 2 c\,(16 a c e - 8 a d^2 - 8 b^2 e + 4 b c d - c^3)\\
   &\qquad{} - 4 a e\,(2\Sigma\alpha)^2 - 4 a c^2 e\\
   &\qquad{} + 4 a e\,(2\Sigma\alpha\beta) + 4 a e\,(-\,16 a e + 4 b d - c^2)\\
   &\qquad{} + 48 a^2 e^2 + 48 a^2 e^2.
\end{aligned}
\]
\[
\begin{aligned}
   \text{Coeff.\ }\phi^0 &= -\,(16 a c e - 8 a d^2 - 8 b^2 e + 4 b c d - c^3)^2\\
   &\qquad{} + 4 a c e\,(16 a c e - 8 a d^2 - 8 b^2 e + 4 b c d - c^3)\\
   &\qquad{} - \alpha^2\beta^2\gamma^2 - 16 a^2 e^2\,(-\,16 a e + 4 b d - c^2)\\
   &\qquad{} + 4 a e\,\alpha\beta\gamma\Sigma\alpha + 64 a^3 e^3\\
   &\qquad{} - 16 a^2 e^2\,\Sigma\alpha\beta\\
   &\qquad{} + 64 a^3 e^3.
\end{aligned}
\]

\newpage
\setcounter{page}{318}

\noindent\textbf{318}\hfill\textsc{on a new auxiliary equation in}\hfill\textbf{[268}

\medskip

\noindent 21. \ Effecting the developments, these are
\[
\begin{array}{rr}
a e\;\;\quad -\,20 & a^2 b f\quad +\,240\\
b d\;\;\quad +\,8 & a^2 c e f\quad -\,128\\
c^2\;\;\quad -\,3 & a^2 d^4\quad -\,64 \\
& a b c d\quad +\,240\\
& a b d e\quad -\,128\\
& a c^2 e\quad -\,128\\
& a c d^2\quad +\,16\\
& b^2 c e\quad +\,16\\
& b^2 d^2\quad +\,16\\
& b c^2 d\quad +\,16\\
& c^4\quad +\,3\\
\end{array}
\]
the first of which is in fact complete; the others being completed, we obtain the equation in $\phi$, viz.:

\medskip

\noindent 22. \ For the denumerate form $(a, b, c, d, e, f\Cdel v,\,1)^5 = 0$, the equation in $\phi$ is
\[
   \bigl(\,a^6\ X,\ -\,800\sqrt{5}\,a^5\sqrt{D}\,X,\ a^4\ X,\ a^4\ X,\ a^2\ X,\ X\,\bigr)(\phi,\,1)^6 = 0,
\]
[the explicit columns of monomial-coefficients given in the original printed page~318: column $a^4$ produces the developed table including $a e\,(-20)$, $b d\,(+8)$, $c^2\,(-3)$, etc., column $a^2$ produces the table for the cubic-in-$\phi$ coefficient, and the constant column produces the $\phi^0$ table; these are reproduced verbatim in the printed text and amount to several score of monomials each --- they are taken as given by Cayley without alteration], where $D\,(= a^2 f^2 + \&c.)$ denotes the discriminant for the denumerate form.

\newpage
\setcounter{page}{319}

\noindent\textbf{268]}\hfill\textsc{the theory of equations of the fifth order.}\hfill\textbf{319}

\medskip

\noindent I proceed now to form the expression for $\phi_1\phi_4 + \phi_2\phi_5 + \phi_3\phi_6$.

\medskip

\noindent 23. \ Writing for convenience $x$ in the place of the root 5, we have
\[
\begin{aligned}
   \phi_1 &= A - B + \{23 - 14 + x\,(1 + 4 - 2 - 3)\},\\
   \phi_4 &= A - B - \{23 - 14 + x\,(1 + 4 - 2 - 3)\},
\end{aligned}
\]
or
\[
   \phi_1\phi_4 = (A - B)^2 - \{23 - 14 + x\,(1 + 4 - 2 - 3)\}^2.
\]
The terms without $x$ are, as before, $(A - B)^2 - C^2 + 4\cdot\overline{1234}$, or $-\alpha\beta + 4\cdot\overline{1234}$, and we have
\[
\begin{aligned}
   \phi_1\phi_4 &= -\,\alpha\beta + 4\cdot 12\cdot 34\\
   &\quad{} + 2 x\,(1 + 4 - 2 - 3)(14 - 23)\\
   &\quad{} - x^2\,(1 + 4 - 2 - 3)^2;
\end{aligned}
\]
and in like manner
\[
\begin{aligned}
   \phi_2\phi_5 &= -\,\beta\gamma + 4\cdot\overline{1234}\\
   &\quad{} + 2 x\,(1 + 2 - 3 - 4)(12 - 34)\\
   &\quad{} - x^2\,(1 + 2 - 3 - 4)^2,
\end{aligned}
\]
and
\[
\begin{aligned}
   \phi_3\phi_6 &= -\,\gamma^2 + 4\cdot 12\cdot 34\\
   &\quad{} + 2 x\,(1 + 3 - 4 - 2)(13 - 42)\\
   &\quad{} - x^2\,(1 + 3 - 4 - 2)^2.
\end{aligned}
\]

\medskip

\noindent 24. \ The roots $1, 2, 3, 4$, contained in these expressions explicitly, and in $\alpha, \beta, \gamma$, are the roots of the equation $\dfrac{1}{v - x}\,(a, b, c, d, e, f\Cdel v,\,1)^5 = 0$, or, what is the same thing,
\[
   (a', b', c', d', e'\Cdel v,\,1)^4 = 0,
\]
where
\[
\begin{aligned}
   a' &= a,\\
   b' &= a x + b,\\
   c' &= a x^2 + b x + c,\\
   d' &= a x^3 + b x^2 + c x + d,\\
   e' &= a x^4 + b x^3 + c x^2 + d x + e.
\end{aligned}
\]

\medskip

\noindent Omitting, as before, a power of $a$, which is ultimately restored, we have
\[
\begin{aligned}
   \phi_1\phi_4 + \phi_2\phi_5 + \phi_3\phi_6 &= -\,\Sigma\alpha\beta + 12 a' e'\\
   &\quad{} + 2 x\,\Sigma\,(1 + 4 - 2 - 3)(14 - 23)\\
   &\quad{} - x^2\,\Sigma\,(1 + 4 - 2 - 3)^2,
\end{aligned}
\]
where the $\Sigma$'s in the second and third lines denote each of them the sum of the three terms obtained by the cyclical permutations of $2, 3, 4$.

\newpage
\setcounter{page}{320}

\noindent\textbf{320}\hfill\textsc{on a new auxiliary equation in}\hfill\textbf{[268}

\medskip

\noindent The first line is
\[
\begin{aligned}
   &(16 a' e' - 4 b' d' + c'^2) + 12 a' e'\\
   &\quad{} = 28 a' e' - 4 b' d' + c'^2.
\end{aligned}
\]
The second line is $2 x$ into $2\,\Sigma 1\cdot 2 - 3\,\Sigma 123$,
or it is
\[
\begin{aligned}
   &= 2 x\,[(-\,b' c' + 3 a' d') + 3 a' d']\\
   &= 2 x\,(6 a' d' - 2 b' c');
\end{aligned}
\]
and the third line is $-\,x^2$ into $3\,\Sigma 1^2 - 2\,\Sigma 12$,
or it is
\[
\begin{aligned}
   &= 3\,(b'^2 - 2 a' c') - 2 a' c'\\
   &= x^2\,(8 a' c' - 3 b'^2).
\end{aligned}
\]
Hence, combining the three terms,
\[
\begin{aligned}
   \phi_1\phi_4 + \phi_2\phi_5 + \phi_3\phi_6 &= 28 a' e' - 4 b' d' + c'^2\\
   &\quad{} + x\,(12 a' d' - 2 b' c')\\
   &\quad{} + x^2\,(8 a' c' - 3 b'^2),
\end{aligned}
\]
or substituting for $(a', b', c', d', e')$ their values, the right-hand side is
\[
   = (40 a^2,\ 32 a b,\ 28 a c - 8 b^2,\ 44 a d - 8 b c,\ 28 a e - 4 b d + c^2\Cdel x,\,1)^4,
\]
where $x$ stands for $x_5$, and on the left-hand side the factor $a^2$ is to be restored.

\medskip

\noindent 25. \ Writing for shortness
\[
   (*\Cdel x,\,1)^4 = (40 a^2,\ 32 a b,\ 28 a c - 8 b^2,\ 44 a d - 8 b c,\ 28 a e - 4 b d + c^2\Cdel x,\,1)^4,
\]
the equation is
\[
   a^2\,(\phi_1\phi_4 + \phi_2\phi_5 + \phi_3\phi_6) = (*\Cdel x_5,\,1)^4;
\]
and the system of equations to which this belongs is
\[
\begin{aligned}
   a^2\,(\phi_1\phi_2 + \phi_3\phi_4 + \phi_5\phi_6) &= (*\Cdel x_1,\,1)^4,\\
   a^2\,(\phi_1\phi_5 + \phi_2\phi_3 + \phi_4\phi_6) &= (*\Cdel x_2,\,1)^4,\\
   a^2\,(\phi_1\phi_3 + \phi_2\phi_6 + \phi_3\phi_5) &= (*\Cdel x_3,\,1)^4,\\
   a^2\,(\phi_2\phi_4 + \phi_3\phi_6 + \phi_4\phi_5) &= (*\Cdel x_4,\,1)^4,\\
   a^2\,(\phi_1\phi_4 + \phi_2\phi_5 + \phi_3\phi_6) &= (*\Cdel x_5,\,1)^4;
\end{aligned}
\]
so that the roots $x_1, x_2, x_3, x_4, x_5$ will be rational functions of the combinations $\phi_1\phi_2 + \phi_3\phi_4 + \phi_5\phi_6$, \&c.\ respectively, of the equation in $\phi$.

\newpage
\setcounter{page}{321}

\noindent\textbf{268]}\hfill\textsc{the theory of equations of the fifth order.}\hfill\textbf{321}

\medskip

\noindent 26. \ Passing now to the standard form $(a, b, c, d, e, f\Cdel v,\,1)^5 = 0$, the equation in $\phi$ is
\[
\bigl(\,a^6,\ -\,100\,a^5\sqrt{5}\sqrt{\square}\,X,\ 2000\,a^4 X,\ a^4 X,\ a^2 X,\ X,\ 40000\,X\,\bigr)(\phi,\,1)^6 = 0,
\]
[where again the column-tables of monomial coefficients are reproduced verbatim from page~321 of the original printed text; these tables enumerate, for each column, the monomials $a e$, $b d$, $c^2$, $a^2 e f$, $a b c f$, $a b d e$, $a c^2 e$, $a c d^2$, $b^2 c e$, $b^2 d^2$, $b c^2 d$, $c^4$, etc., with their integer coefficients $+1, -2, +1, +1, -2, +14, -2, +8, -10, -40, -50, +15$, \&c., the full table being given on the original page],
where $\square\,(= a^2 f^2 + \&c.)$ denotes the Discriminant for the Standard form.

\medskip

\noindent 27. \ And if we put
\[
   (*\Cdel x,\,1)^4 = 20\,(2 a^2,\ 8 a b,\ 22 a c - 10 b^2,\ 18 a d - 10 b c,\ 7 a e - 10 b d + 5 c^2\Cdel x,\,1)^4,
\]
then we have
\[
\begin{aligned}
   a^2\,(\phi_1\phi_2 + \phi_3\phi_4 + \phi_5\phi_6) &= (*\Cdel x_1,\,1)^4,\\
   a^2\,(\phi_1\phi_5 + \phi_2\phi_3 + \phi_4\phi_6) &= (*\Cdel x_2,\,1)^4,
\end{aligned}
\]
\&c., which lead to rational expressions for the roots $x_1, x_2, x_3, x_4, x_5$, in terms of the combinations $\phi_1\phi_2 + \phi_3\phi_4 + \phi_5\phi_6$, \&c.\ respectively.

\hfill C.~IV.\hfill 41

\newpage
\setcounter{page}{322}

\noindent\textbf{322}\hfill\textsc{on a new auxiliary equation in}\hfill\textbf{[268}

\medskip

\noindent 28. \ Consider now the quintic function
\[
   U = (a, b, c, d, e, f\Cdel x, y)^5 = a\,(x - \alpha y)(x - \beta y)(x - \gamma y)(x - \delta y)(x - \epsilon y);
\]
and treating the numbers $1, 2, 3, 4, 5$ as corresponding to $\alpha, \beta, \gamma, \delta, \epsilon$ respectively, write
\[
   \Phi = \widetilde{\overline{12345}} - \widetilde{\overline{24135}},\ \&c.,
\]
where
\[
   \widetilde{\overline{12345}} = \widetilde{12} + \widetilde{23} + \widetilde{34} + \widetilde{45} + \widetilde{51},
\]
in which $\widetilde{12}$, \&c.\ denote respectively
\[
   \frac{1}{y^2}\cdot\frac{1}{x - \alpha y}\cdot\frac{1}{x - \beta y},\ \&c.
\]
Then we have
\[
   \Phi = \frac{a}{U}\,\{\phi x - \chi y\},
\]
where
\[
   \phi = \overline{12345} - \overline{24135},
\]
and
\[
   \overline{12345} = 12 + 23 + 34 + 45 + 51 = \alpha\beta + \beta\gamma + \gamma\delta + \delta\epsilon + \epsilon\alpha,
\]
and
\[
   \overline{24135} = 24 + 41 + 13 + 35 + 52 = \beta\delta + \delta\alpha + \alpha\gamma + \gamma\epsilon + \epsilon\beta,
\]
and
\[
   \chi = (\overline{12345}) - (\overline{24135}),
\]
\[
   (\overline{12345}) = 123 + 234 + 345 + 451 + 512 = \alpha\beta\gamma + \beta\gamma\delta + \gamma\delta\epsilon + \delta\epsilon\alpha + \epsilon\alpha\beta,
\]
\[
   (\overline{24135}) = 241 + 413 + 135 + 352 + 524 = \beta\delta\alpha + \delta\alpha\gamma + \alpha\gamma\epsilon + \gamma\epsilon\beta + \epsilon\beta\delta.
\]

\medskip

\noindent 29. \ In fact,
\[
   \Phi = \frac{a}{U}\,\{(\overline{12345}) - (\overline{24135})\},
\]
where
\[
   (\overline{12345}) = 123 + 234 + 345 + 451 + 512,
\]
\[
   (\overline{24135}) = 241 + 413 + 135 + 352 + 524,
\]
where 123, \&c.\ denote respectively
\[
   \frac{1}{y^3}\,(x - \alpha y)(x - \beta y)(x - \gamma y),\ \&c.,
\]
and $(\overline{12345}) - (\overline{24135})$ thus presents itself as a cubic function divided by $y^3$. But in this cubic function the coefficients of $x^3$, $x^2 y$ vanish. For the coefficient of any power of $x$ will be
\[
   123 + 234 + 345 + 451 + 512 - 241 - 413 - 135 - 352 - 524,
\]

\newpage
\setcounter{page}{323}

\noindent\textbf{268]}\hfill\textsc{the theory of equations of the fifth order.}\hfill\textbf{323}

\medskip

\noindent where, first, for $x^3$, 123, \&c.\ denote respectively unity; the coefficient of $x^3$ therefore vanishes. Next, for $x^2 y$, 123, \&c.\ denote respectively $-\,(1 + 2 + 3)$, \&c.\ $(= \alpha + \beta + \gamma)$, and the coefficient of $x^2 y$ also vanishes. But for $x y^2$, 123, \&c.\ denote respectively $12 + 23 + 31\,(= \alpha\beta + \beta\gamma + \gamma\alpha)$, \&c.\ respectively; the positive terms are
\[
   (12 + 23 + 31) + (23 + 34 + 42) + (34 + 45 + 53) + (45 + 51 + 14) + (51 + 12 + 25),
\]
which are
\[
\begin{aligned}
   &= 2\,(12 + 23 + 34 + 45 + 51) + (24 + 41 + 13 + 35 + 52)\\
   &= 2\cdot\overline{12345} + \overline{24135};
\end{aligned}
\]
and the negative terms, taken positively, are
\[
   (24 + 41 + 12) + (41 + 13 + 34) + (13 + 35 + 51) + (35 + 52 + 23) + (52 + 24 + 45),
\]
which are
\[
\begin{aligned}
   &= (12 + 23 + 34 + 45 + 51) + 2\,(24 + 41 + 13 + 35 + 52)\\
   &= \overline{12345} + 2\cdot\overline{24135};
\end{aligned}
\]
so that the difference, or coefficient of $x y^2$, is
\[
   = \overline{12345} - \overline{24135},
\]
which is $= \phi$.

\medskip

\noindent And for $y^3$, 123, \&c.\ denote respectively $-\,123\,(= -\,\alpha\beta\gamma)$, \&c., so that the coefficient of $y^3$ is $= \chi$.

\medskip

\noindent 30. \ The cubic function is therefore $= \phi x y^2 - \chi y^3$; and dividing by $y^2$, we have
\[
   \Phi = \frac{a}{U}\,(\phi x - \chi y).
\]
$\Phi$ is thus a fractional covariantive function, the leading coefficient whereof is $\phi$, and the equation for the determination of $\Phi$ is consequently that deduced from the equation for $\phi$, by replacing therein the seminvariants by the corresponding covariants.

\medskip

\noindent The equation [denoting the covariants as in 141 and 143, $A$ the quintic itself, \&c.] is
\[
\begin{aligned}
   &-\,100 A^5 B,\\
   &0,\\
   &+\,2000 A^4\,(6 B^2 - 4 H)\\
   &-\,800\,A^5\sqrt{B}\sqrt{\text{disct.}},\\
   &A J - 25 D^2,\quad\Cdel\Phi,\,1)^6 = 0
\end{aligned}
\]
where the coefficients are of the orders 30, $-$, 22, $-$, 14, 10, 6\\
\hfill 41--2

\newpage
\setcounter{page}{324}

\noindent\textbf{324}\hfill\textsc{on a new auxiliary equation \&c.}\hfill\textbf{[269}

\medskip

\noindent respectively. The last coefficient [now given in the form $A J - 25 D^2$] being of the degree 6 in the coefficients $(a, b, c, d, e, f)$, is not given in the Tables; it is therefore merely indicated by $(\mathfrak{A},\mathfrak{B},\mathfrak{C},\mathfrak{F},\mathfrak{G},\mathfrak{H}\Cdel x, y)^6$, the leading coefficient $\mathfrak{A}$ being of course the last coefficient in the equation for $\phi$, to the standard form.

\medskip

\noindent I refrain from at present entering into the consideration of the values of the expressions $\Phi_1\Phi_2 + \Phi_3\Phi_4 + \Phi_5\Phi_6$, \&c.

\medskip

\begin{center}\textit{Addition}, Nov.~10, 1862.\ [Originally printed in a later memoir ``On Tschirnhausen's Transformation'' post, 275.]\end{center}

\medskip

\noindent I take the opportunity of remarking, with reference to [the foregoing] memoir, that I recently discovered that the auxiliary equation there considered is in fact due to Jacobi, who, in his paper, ``Observatiunculae ad theoriam aequationum pertinentes,'' \textit{Crelle}, t.~\textsc{xiii}.\ (1835), pp.~340--352, under the heading ``Observatio de aequatione sexti gradus ad quam aequationes quinti gradus revocari possunt,'' gives the theory, and observes that the equation is of the form
\[
   \phi^6 + a_2\phi^4 + a_4\phi^2 + a_6 = 32\,\sqrt{D}\,\phi,
\]
and mentions that the value of $a_2$ is easily found to be (I adapt his notation to the denumerate form $(a, b, c, d, e, f\Cdel v,\,1)^5 = 0$)
\[
   = 40 a e - 16 b d + 6 c^2
\]
(this ought, however, to be divided by $-\,2$), but that the values of $a_4$, $a_6$ ``paullo ampliores calculos poscunt.''

\medskip

\noindent The value of the coefficient in question is correctly obtained (page 270 of my memoir [as printed in the Philosophical Transactions]) in the form
\[
\begin{aligned}
   & +\,2\,(-\,16 a e + 4 b d - c^2)\\
   & +\,12 a e;
\end{aligned}
\]
but the reduced value is given in two places (page 271) as equal to
\[
\begin{aligned}
   -\,32\,a e, &\quad\text{this should be } -\,20\,a e,\\
   +\,8\,d b, &\quad{} \quad\quad\quad\quad\quad +\,8\,b d,\\
   -\,3\,c^2, &\quad{} \quad\quad\quad\quad\quad +\,3\,c^2.
\end{aligned}
\]
The last-mentioned correct value was used in obtaining the coefficient for the standard form, which coefficient is given correctly, page 274. [The correction here indicated $-\,20$ in place of $-\,32$, is made \textit{ante} p.~318.]

\newpage
\setcounter{page}{325}

% ============================================================
% PAPER 269
% A SEVENTH MEMOIR ON QUANTICS
% ============================================================

\begin{center}
\textbf{269.}\\[1em]
{\Large A SEVENTH MEMOIR ON QUANTICS.}\\[1em]
\textit{[From the Philosophical Transactions of the Royal Society of London, vol.~\textsc{cli}.\ (for the year 1861), pp.~277--292.\ \ Received February 28,---Read March 14, 1861.]}
\end{center}

\medskip

\noindent\textsc{The} present memoir relates chiefly to the theory of ternary cubics. Since the date of my Third Memoir on Quantics, [144], M.~Aronhold has published the continuation of his researches on ternary cubics, in the memoir ``Theorie der homogenen Functionen dritten Grades von drei Ver\"anderlichen,'' \textit{Crelle}, t.~\textsc{lv}.\ pp.~97--191 (1858). He there considers two derived contravariants, linear functions of the fundamental ones, and which occupy therein the position which the fundamental contravariants $PU$, $QU$ do in my Third Memoir; in the notation of the present memoir these derived contravariants are
\[
\begin{aligned}
   Y U &= 3 T\cdot P U - 4 S\cdot Q U,\\
   Z U &= -\,48 S^2\cdot P U + T\cdot Q U;
\end{aligned}
\]
and for the canonical form $x^3 + y^3 + z^3 + 6 l x y z$, they acquire respectively the factor $(1 + 8 l^3)^2$, viz.\ in this case
\[
\begin{aligned}
   Y U &= (1 + 8 l^3)^2\,\{l\,(\xi^3 + \eta^3 + \zeta^3) - 3 l^2\,\xi\eta\zeta\},\\
   Z U &= (1 + 8 l^3)^2\,\{(1 + 2 l^3)(\xi^3 + \eta^3 + \zeta^3) + 18\,l^2\,\xi\eta\zeta\}.
\end{aligned}
\]
The derived contravariants have with the covariants $U$, $H U$, even a more intimate connexion than have the contravariants $P U$, $Q U$; and the advantage of the employment of $Y U$, $Z U$ fully appears by M.~Aronhold's memoir.

\medskip

\noindent But the conclusion is, not that the contravariants $P U$, $Q U$ are to be rejected, but that the system is to be completed by the addition thereto of two derived covariants, linear functions of $U$, $H U$; these derived covariants, suggested to me by M.~Aronhold's memoir, are in the present memoir called $C U$, $D U$; their values are
\[
\begin{aligned}
   C U &= -\,T\cdot U + 24 S\cdot H U,\\
   D U &= 8 S^2\cdot U - 3 T\cdot H U;
\end{aligned}
\]

\newpage
\setcounter{page}{326}

\noindent\textbf{326}\hfill\textsc{a seventh memoir on quantics.}\hfill\textbf{[269}

\medskip

\noindent and for the canonical form $x^3 + y^3 + z^3 + 6 l x y z$, they acquire respectively, not indeed $(1 + 8 l^3)^2$, but the simple power $(1 + 8 l^3)$, as a factor, viz.\ in this case
\[
\begin{aligned}
   C U &= (1 + 8 l^3)\,\{(-1 + 4 l^3)(x^3 + y^3 + z^3) + 18 l^2 x y z\},\\
   D U &= (1 + 8 l^3)\,\{l^2\,(5 + 4 l^3)(x^3 + y^3 + z^3) + 3\,(1 - 10 l^3)\,x y z\};
\end{aligned}
\]
it was in fact by means of this condition as to the factor $(1 + 8 l^3)$, that the foregoing expressions for $C U$, $D U$ were obtained.\footnote{M.~Aronhold, in a letter dated Berlin, 17 June 1861, has pointed out to me that the covariants $C U$, $D U$ are in his notation $P_{\Theta}$, $P_{\Psi}$, and that they belong to the forms called Conjugate Forms, \S~27 of his memoir. But the explicit development of the properties of these covariants is not on this account the less interesting. Added 20 Sept.\ 1861.---A.C.}

\medskip

\noindent The formul\ae{} of my Third Memoir and those of M.~Aronhold are by this means brought into harmony and made parts of a whole; instead of the two intermediates
\[
   \alpha U + 6\beta H U,\quad 6\alpha P U + \beta Q U,
\]
in Tables 68 and 69 of my Third Memoir, or of the intermediates
\[
   \alpha U + 6\beta H U,\quad -\,2\alpha Y U + 2\beta Z U,
\]
of M.~Aronhold's theory, we have the four intermediates
\[
   \alpha U + 6\beta H U,\quad -\,2\alpha Y U + 2\beta Z U,\quad 2\alpha C U - 2\beta D U,\quad 6\alpha P U + \beta Q U,
\]
in Tables 74, 75, 76, and 77 of the present memoir. These four Tables embrace the former results, and the new ones which relate to the covariants $C U$, $D U$; and they are what is most important in the present memoir. I have, however, excluded from the Tables, and I do not in the memoir consider (otherwise than incidentally) the covariant of the sixth order $\Theta U$, or the contravariant (reciprocant) $F U$.

\medskip

\noindent I have given in the memoir a comparison of my notation with that of M.~Aronhold. A short part of the memoir relates to the binary cubic and the binary quartic, viz.\ each of these quantics has a covariant of its own order, forming with it an intermediate $\alpha U + \beta W$, the covariants whereof contain quantics in $(\alpha, \beta)$, the coefficients of which are invariants of the original quantic. The formul\ae{} which relate to these cases are in fact given in my Fifth Memoir, [156], but they are reproduced here in order to show the relations between the quantics in $(\alpha, \beta)$ contained in the formul\ae. As regards the binary quartic, these results are required for the discussion of the like question in regard to the ternary cubic, viz.\ that of finding the relations between the different quantics in $(\alpha, \beta)$ contained in the formul\ae{} relating to the ternary cubic. Some of these relations have been obtained by M.~Hermite in the memoir ``Sur les formes cubiques \`a trois ind\'etermin\'ees,'' \textit{Liouville}, t.~\textsc{iii}.\ pp.~37--40 (1858), and in that ``Sur la R\'esolution des Equations du quatri\`eme degr\'e,'' \textit{Comptes Rendus}, \textsc{xlvi}.\ p.~715 (1858), and by M.~Aronhold in his memoir already referred to; and in particular I reproduce and demonstrate some of the results in the last-mentioned memoir of M.~Hermite. But the relations in question are in the present memoir exhibited in a more complete and systematic form.

\newpage
\setcounter{page}{327}

\noindent\textbf{269]}\hfill\textsc{a seventh memoir on quantics.}\hfill\textbf{327}

\medskip

\noindent The paragraphs and Tables of the present memoir are numbered consecutively with those of my former memoirs on Quantics.

\medskip

\noindent 231. \ For the binary cubic $(a, b, c, d\Cdel x, y)^3$, if $U$ be the cubic itself, $H U$ the Hessian, $\Phi U$ the cubicovariant, and $D$ the discriminant (see Fifth Memoir, Nos.~115, 118), then
\begin{center}Covariant and other Tables, No.~71.\end{center}
\[
   D\,(\alpha U + \beta\Phi U) = (\alpha^2 - \beta^2 D)\,U,
\]
so that the quantics in $(\alpha, \beta)$ all of them depend on $\alpha^2 - \beta^2 D$.

\medskip

\noindent 232. \ For the binary quartic $(a, b, c, d, e\Cdel x, y)^4$, if $U$ be the quartic itself, $H U$ the Hessian, $\Phi U$ the cubicovariant, $I, J$, the quadrinvariant and the cubinvariant, and $D\,(= I^3 - 27 J^2)$ the discriminant (see Fifth Memoir, Nos.~128, 134), then
\begin{center}Table No.~72.\end{center}
\[
   \Phi\,(\alpha U + 6\beta H U) = (1,\,0,\,-9 I,\,-54 J\Cdel\alpha,\,\beta)^3\cdot U,
\]
\[
\begin{aligned}
   H\,(\alpha U + 6\beta H U) &= -\tfrac{1}{12}\tfrac{d}{d\beta}\,(1,\,0,\,-9 I,\,-54 J\Cdel\alpha,\,\beta)^3\cdot U\\
   &\quad{} + \tfrac{1}{3}\tfrac{d}{d\alpha}\,(1,\,0,\,-9 I,\,-54 J\Cdel\alpha,\,\beta)^3\cdot H U,
\end{aligned}
\]
\[
   I\,(\alpha U + 6\beta H U) = (I,\,18 J,\,3 I^2\Cdel\alpha,\,\beta)^2,
\]
\[
   J\,(\alpha U + 6\beta H U) = (J,\,I^2,\,9 I J,\,-I^3 + 54 J^2\Cdel\alpha,\,\beta)^3,
\]
\[
   D\,(\alpha U + 6\beta H U) = (1,\,0,\,-18 I,\,-108 J,\,81 I^2,\,972 I J,\,2916 J^2\Cdel\alpha,\,\beta)^6\cdot D\footnote{The coefficient $2916 J^2$ is in the Fifth Memoir erroneously given as $-\,2916 J^2$. [This correction should have been made, vol.~ii.\ p.~549.]}
\]
\[
   = [(1,\,0,\,-9 I,\,-54 J\Cdel\alpha,\,\beta)^3]^2\cdot D.
\]

\medskip

\noindent 233. \ Writing for the moment
\[
   G = (1,\,0,\,-9 I,\,-54 J\Cdel\alpha,\,\beta)^3,
\]
then the Hessian, cubicovariant, and discriminant of this cubic function of $(\alpha, \beta)$ are respectively
\[
\begin{aligned}
   H G &= -\,3\,(I,\,18 J,\,3 I^2\Cdel\alpha,\,\beta)^2,\\
   \Phi G &= 54\,(J,\,I^2,\,9 I J,\,-I^3 + 54 J^2\Cdel\alpha,\,\beta)^3,\\
   D G &= -\,108 D;
\end{aligned}
\]
so that the covariants of the intermediate $\alpha U + 6\beta H U$ are all of them expressible by means of the cubic function $G$.

\newpage
\setcounter{page}{328}

\noindent\textbf{328}\hfill\textsc{a seventh memoir on quantics.}\hfill\textbf{[269}

\medskip

\noindent It may be noticed that $G$ is what the left-hand side of the equation
\[
   I\,(H U)^2 - J\cdot H U\cdot U^2 + J\,I^3 = -\,(\Phi U)^2
\]
(see Fifth Memoir, No.~128) becomes on writing therein $\alpha, -\,6\beta$, for $U, H U$ respectively, and throwing out the factor 4.

\medskip

\noindent 234. \ I take the opportunity of remarking with respect to a binary quartic $U = (a, b, c, d\Cdel x, y)^4$, that the Hessian of the cubicovariant, to fix the numerical factor, say $\tfrac{1}{144}\,(\partial_x^2\,\Phi U\cdot\partial_y^2\,\Phi U - (\partial_x\partial_y\,\Phi U)^2)$, is
\[
   I^3 U^2 - 36 J\cdot U\cdot H U + 12 I\,(H U)^2,
\]
which is
\[
   = (I U - 18 J\,H U)^2 + 4\,(I^3 - 27 J^2)\,(H U)^2;
\]
or if $I^3 - 27 J^2 = 0$, that is if the quartic has a pair of equal factors, the Hessian of the cubicovariant is a perfect square.

\medskip

\noindent 235. \ Coming now to the ternary cubic $U = (a, b, c, f, g, h, i, j, k, l\Cdel x, y, z)^3$, I give in the first place the following comparison of my notation with that of M.~Aronhold.
\[
\begin{array}{ll}
\textbf{Aronhold.} & \textbf{Cayley.}\\[3pt]
f & -CHU\\
\Delta f & -\,6 H U\\
S & 4 S\\
T & -\,T\\
A_8\,/ & R\\
T_f & 6 P U\\
A_5\,/ & -\,2 Y U\\
J\quad & 2 Z U\\
\quad & C U\\
\quad & D U\\
0 & \\
H & -\,F U\\
F\quad & 2\,(\Theta_w U - T U^2 + 4 S U\cdot H U),
\end{array}
\]
where the notations $Y U$, $Z U$ (to correspond to M.~Aronhold's $P_f$, $Q_f$) and the notations $C U$, $D U$ are first employed in the present memoir. I remark in regard to $P_f\,(= -\,2 Y U)$, where, as already mentioned,
\[
   Y U = (1 + 8 l^3)^2\,\{l\,(\xi^3 + \eta^3 + \zeta^3) - 3 l^2\,\xi\eta\zeta\},
\]

\newpage
\setcounter{page}{329}

\noindent\textbf{269]}\hfill\textsc{a seventh memoir on quantics.}\hfill\textbf{329}

\medskip

\noindent that in my Memoir on Curves of the Third Order (\textit{Phil.\ Trans.}, t.~\textsc{cxlvii}.\ (1857), see p.~427), [146], I was led incidentally to the curve
\[
   l\,(\xi^3 + \eta^3 + \zeta^3) - 3 l^2\,\xi\eta\zeta = 0,
\]
and that I there gave the equation
\[
   3 T\cdot P U - 4 S\cdot Q U = (1 + 8 l^3)^2\,\{l\,(\xi^3 + \eta^3 + \zeta^3) - 3 l^2\,\xi\eta\zeta\}.
\]
But the curve
\[
   (1 + 2 l^3)\,(\xi^3 + \eta^3 + \zeta^3) + 18 l^2\,\xi\eta\zeta = 0,
\]
which corresponds to $(Q_f = 2 Z U)$, does not occur in that memoir.

\medskip

\noindent 236. \ I remark, further, in regard to M.~Aronhold's $\Theta, H$, that these are what he calls ``Zwischenformen,'' viz.\ they are covariants of the cubic and of the adjoint linear form $\xi x + \eta y + \zeta z$, or as they might be termed Contracovariants. For the canonical form $U = x^3 + y^3 + z^3 + 6 l x y z$, the value of $\Theta$ is
\[
   (y z - l^2 x^2,\ z x - l^2 y^2,\ x y - l^2 z^2,\ l^2 y z - l x^2,\ l^2 z x - l y^2,\ l^2 x y - l z^2\Cdel\xi, \eta, \zeta)^2,
\]
which is a form which occurs incidentally in my memoir last referred to (see p.~427). The value of $H$ in the same case is
\[
   (-\,2 l\,(1 + 2 l^3)\,x^2 - 6 l y z,\ldots,-(1 + 4 l^3)\,x^2 + 2 l\,(1 + 2 l^3)\,y z,\ldots\Cdel\xi, \eta, \zeta)^2,
\]
which does not occur in that memoir. In my Third Memoir on Quantics I purposely abstained from the consideration of any such forms.

\medskip

\noindent 237. \ My covariants $\Theta U$ and $\psi U$ involved unsymmetrically the cubic and its Hessian, and it did not occur to me how a similar covariant, such as M.~Aronhold's $\psi$, which involves the two functions symmetrically, was to be formed. Let $(A, B, C)$ be the first derived functions, $(a, b, c, f, g, h)$ the second derived functions of the cubic $U$, and $(A', B', C')$ the first derived functions, $(a', b', c', f', g', h')$ the second derived functions of the Hessian $H U$, then disregarding numerical factors, we have
\[
   \Theta U = (b c - f^2,\ldots,g h - a f,\ldots\Cdel A', B', C')^2,
\]
and
\[
   \Theta_1 U = (b' c' - f'^2,\ldots,g' h' - a' f',\ldots\Cdel A, B, C)^2,
\]
\[
   \psi = (b c' + b' c - 2 f f',\ldots,g h' + g' h - a f' - a' f,\ldots\Cdel A, B, C\Cdel A', B', C');
\]
and considering $U = 0$ as the equation of a curve of the third order, the equations $\Theta U = 0$, $\Theta_1 U = 0$, $\psi = 0$ have the following significations, viz.\ $\Theta U = 0$ is the locus of a point, such that its second or line polar with respect to the Hessian touches its first or conic polar with respect to the cubic: $\Theta_1 U$ is the locus of a point such that its second or line polar with respect to the cubic touches its first or conic polar with respect to the Hessian: and $\psi = 0$ is the locus of a point such that its second or line polar with respect to the cubic, and its second or line polar with respect to the Hessian are reciprocals (that is, each passes through the pole of the other of them) with respect to the conic which is the envelope of a line cutting the first or conic polar of the point with respect to the cubic, and the first or conic polar of the point with respect to the Hessian, in two pairs of points which are harmonically related to each other: such being in fact the immediate interpretation of the analytical formula. But this in passing.

\hfill C.~IV.\hfill 42

\newpage
\setcounter{page}{330}

\noindent\textbf{330}\hfill\textsc{a seventh memoir on quantics.}\hfill\textbf{[269}

\medskip

\noindent 238. \ The formul\ae{} (Tables 68 and 69 of my Third Memoir) for the discriminants of the intermediates $\alpha U + 6\beta H U$ and $6\alpha P U + \beta Q U$ respectively are
\[
\begin{aligned}
   R\,(\alpha U + 6\beta H U) &= [(1,\,0,\,-24 S,\,-8 T,\,-48 S^2\Cdel\alpha,\,\beta)^4]^3\,R,\\
   R\,(6\alpha P U + \beta Q U) &= [(48 S,\,8 T,\,-96 S^2,\,-24 T S,\,-T^2 - 16 S^3\Cdel\alpha,\,\beta)^4]^3\,R.
\end{aligned}
\]

\medskip

\noindent In M.~Hermite's paper in the \textit{Comptes Rendus}, already referred to, there are given between these quantics in $(\alpha, \beta)$ certain relations which (although less simple than the relations that will afterwards be obtained) I now proceed to investigate. Putting in the first formula $\alpha\div\beta = p$, and in the second formula $\alpha\div\beta = \theta$, we have
\[
\begin{aligned}
   R\,(p U + 6 H U) &= 0,\quad\text{if } (1,\,0,\,-24 S,\,-8 T,\,-48 S^2\Cdel p,\,1)^4 = 0,\\
   R\,(6\theta P U + Q U) &= 0,\quad\text{if } (48 S,\,8 T,\,-96 S^2,\,-24 T S,\,-T^2 - 16 S^3\Cdel\theta,\,1)^4 = 0,
\end{aligned}
\]
which equations in $p, \theta$, are about to be considered in place of the quantics from which they respectively arise.

\medskip

\noindent 239. \ It is convenient to write\footnote{A slight inaccuracy in three of Hermite's formul\ae, which should be $\beta$, $B$, is here corrected. Hermite's $S_3$, $B$ is Hermite's $S_4$, and $p, q, \theta, \Delta$ are Hermite's $J, S_1, \Delta, \theta$: there is corresponding to formula in the present memoir.}
\[
   B = \tfrac{1}{64}\,T^2 - 64 S^3
\]
(so that $T^2 = A^2 + B^2$ and, for the canonical form,
\[
   A = -\,4 l + 4 l^4,\quad B = 1 + 8 l^3).
\]
Making this change, and joining to the equation in $p$ that derived from it by writing $q$ for $p$, and interchanging $A, B$, we have the three equations
\[
\begin{aligned}
   (1,\,0,\,-6 A,\,-8 T,\,-3 A^2\Cdel p,\,1)^4 &= 0,\\
   (1,\,0,\,-6 B,\,-8 T,\,-3 B^2\Cdel q,\,1)^4 &= 0,\\
   (12 A,\,8 T,\,-6 A^2,\,-6 T A,\,-T^2 - 7 A^3\Cdel\theta,\,1)^4 &= 0.
\end{aligned}
\]

\medskip

\noindent 240. \ The signification of the equation in $q$ is as follows, viz.\ if the quantic
\[
   U = (*\Cdel x, y, z)^3
\]
is transformed into the canonical form
\[
   X^3 + Y^3 + Z^3 + 6 l X Y Z
\]
by means of the linear equations
\[
   (X, Y, Z) = (\Delta\Cdel x, y, z),
\]

\newpage
\setcounter{page}{331}

\noindent\textbf{269]}\hfill\textsc{a seventh memoir on quantics.}\hfill\textbf{331}

\medskip

\noindent where $\Delta$ is a matrix, then using the same letter $\Delta$ to represent the determinant formed out of this matrix, or determinant of substitution, we have
\[
   q = \frac{3}{\Delta^3},
\]
so that the equation in $q$ is one that presents itself in the question of the reduction of the cubic to its canonical form.

\medskip

\noindent In fact the linear transformation gives
\[
   S\,\Delta^4 = -\,l + l^4,
\]
and thence
\[
   T\,\Delta^6 = 1 - 20 l^3 - 8 l^6,
\]
\[
   (T^2 - 64 S^3)\,\Delta^{12} = (1 + 8 l^3)^2,
\]
which, writing $B^2$ in the place of $T^2 - 64 S^3$, becomes
\[
   B^2\,\Delta^{12} = (1 + 8 l^3)^2,
\]
whence also
\[
   B\,\Delta^6 = 1 + 8 l^3,\quad\text{or}\quad 8 l^3 = B\,\Delta^6 - 1,
\]
\[
\begin{aligned}
   8 T\,\Delta^6 &= 8 - 20\,(B\,\Delta^6 - 1) - (B\,\Delta^6 - 1)^2\\
   &= 27 - 18\,B\,\Delta^6 - B^2\,\Delta^{12}
\end{aligned}
\]
or, as this may be written,
\[
   \frac{81}{\Delta^{12}} - \frac{54 B}{\Delta^6} - \frac{24 T}{\Delta^4} - 3 B^2 = 0,
\]
which, putting therein $q = \dfrac{3}{\Delta^3}$, becomes
\[
   (1,\,0,\,-6 B,\,-8 T,\,-3 B^2\Cdel q,\,1)^4 = 0,
\]
the above-mentioned equation in $q$.

\medskip

\noindent 241. \ The relation between $\theta$ and $q$ is
\[
   \theta = -\,\tfrac{1}{2}\,\frac{T}{A} + \frac{B^2}{2 A q},
\]
as may be verified without difficulty. That between $\theta$ and $p$ is
\[
   \theta = -\,\tfrac{1}{2}\,\left(p - \frac{A}{p}\right),
\]
as appears by the identical equation
\[
\begin{aligned}
&(12 A,\,8 T,\,-6 A^2,\,-6 T A,\,-T^2 - 7 A^3\Cdel\tfrac{1}{2}(p + \tfrac{A}{p}),\,1)^4\\
&\quad = \tfrac{1}{64\,p^4}\,(8 A,\,8 T,\,-64 A^2,\,-72 T A,\,-4(6 A^3 - 64 T^2),\,-72 T A^2,\,-12 A^4,\,8 T A^3,\,8 A^4\Cdel p,\,1)^8\\
&\quad = \tfrac{1}{64\,p^4}\,(1,\,0,\,-6 A,\,-8 T,\,-3 A^2\Cdel p,\,1)^4\cdot(8 A,\,8 T,\,6 A^2,\,0,\,-A^3\Cdel p,\,1)^4
\end{aligned}
\]
\hfill 42--2

\newpage
\setcounter{page}{332}

\noindent\textbf{332}\hfill\textsc{a seventh memoir on quantics.}\hfill\textbf{[269}

\medskip

\noindent where the second factor of the product on the right-hand side is
\[
   -\,\tfrac{A}{p}\cdot 2\,(1,\,0,\,-6 A,\,-8 T,\,-3 A^2\Cdel p,\,1)^4.
\]
The relation between $p$ and $q$ is then at once found to be
\[
   p = \frac{\dfrac{2 A^2}{B}}{q + \dfrac{A}{B} + \dfrac{2 T}{B}},
\]
or (since $p, q$ and $A, B$ may be simultaneously interchanged)
\[
   p = \frac{\dfrac{2 A^2}{B}}{q + \dfrac{B}{q} + \dfrac{2 T}{B}}.
\]

\medskip

\noindent 242. \ Let the equations in $p, q$ be represented by $\phi p = 0$, $\psi q = 0$ respectively; then we have
\[
   \phi p = p^4 - 6 A p^2 - 8 T p - 3 A^2,
\]
and therefore
\[
   \tfrac{1}{4}\phi' p = p^3 - 3 A p - 2 T,
\]
whence
\[
\begin{aligned}
   \tfrac{1}{4}p\phi' p &= p^4 - 3 A p^2 - 2 T p\\
   &= 3 A p^2 + 6 T p + 3 A^2,
\end{aligned}
\]
and therefore
\[
   q = \frac{2 B^2 p}{\tfrac{1}{4}p\phi' p - \phi p} = \frac{24 B^2}{\phi' p},
\]
with a like formula for $p$, that is we have
\[
   q = \frac{24 B^2}{\phi' p},\quad p = \frac{24 A^2}{\psi' q},
\]
which with the equation
\[
   \theta = \tfrac{1}{2}\,\left(p - \frac{A}{p}\right),
\]
are the system of equations connecting $\theta, p, q$.

\medskip

\noindent 243. \ As already remarked, we have to consider the two derived covariants
\[
\begin{aligned}
   C U &= -\,T\cdot U + 24 S\cdot H U,\\
   D U &= 8 S^2\cdot U - 2 T\cdot H U,
\end{aligned}
\]
and the two derived contravariants
\[
\begin{aligned}
   Y U &= 3 T\cdot P U - 4 S\cdot Q U,\\
   Z U &= -\,48 S^2\cdot P U + T\cdot Q U,
\end{aligned}
\]

\newpage
\setcounter{page}{333}

\noindent\textbf{269]}\hfill\textsc{a seventh memoir on quantics.}\hfill\textbf{333}

\medskip

\noindent which for the canonical form $x^3 + y^3 + z^3 + 6 l x y z$ are as follows:
\begin{center}Table No.~70 (addition to).\end{center}
\[
\begin{aligned}
   C U &= (1 + 8 l^3)\,[(-1 + 4 l^3)\,(x^3 + y^3 + z^3) + 18 l^2\,x y z],\\
   D U &= (1 + 8 l^3)\,[l^2\,(5 + 4 l^3)\,(x^3 + y^3 + z^3) + 3\,(1 - 10 l^3)\,x y z],\\
   Y U &= (1 + 8 l^3)^2\,[l\,(\xi^3 + \eta^3 + \zeta^3) - 3 l^2\,\xi\eta\zeta],\\
   Z U &= (1 + 8 l^3)^2\,[(1 + 2 l^3)\,(\xi^3 + \eta^3 + \zeta^3) + 18 l^2\,\xi\eta\zeta].
\end{aligned}
\]

\medskip

\noindent 244. \ We have conversely
\[
\begin{aligned}
   3 R\cdot U &= 3 T\cdot C U + 24 S\cdot D U,\\
   3 R\cdot H U &= 8 S^2\cdot C U + T\cdot D U,\\
   -\,3 R\cdot P U &= T\cdot Y U + 4 S\cdot Q U,\\
   -\,3 R\cdot Q U &= 48 S^2\cdot Y U + 3 T\cdot P U,
\end{aligned}
\]
and also the following formul\ae, viz.\ if
\[
   2\alpha C U - 2\beta D U = \alpha'\,U + 6\beta'\,H U;
\]
which give, conversely,
\[
\begin{aligned}
   \alpha' &= -\,2 T\alpha - 16 S^2\beta,\\
   \beta' &= 8 S\alpha + T\beta,
\end{aligned}
\]
\[
\begin{aligned}
   \alpha &= \tfrac{1}{2 R}\,(T\alpha' + 16 S^2\beta'),\\
   \beta &= \tfrac{1}{2 R}\,(-\,8 S\alpha' - 2 T\beta');
\end{aligned}
\]
and moreover, if
\[
   -\,2\alpha Y U + 2\beta Z U = 6\alpha' P U + \beta' Q U;
\]
which give, conversely,
\[
\begin{aligned}
   \alpha' &= -\,(T\alpha + 16 S^2\beta),\\
   \beta &= -\,(-\,8 S\alpha - 2 T\beta),
\end{aligned}
\]
\[
\begin{aligned}
   \alpha &= -\,\tfrac{1}{2 R}\,(-\,2 T\alpha' - 16 S^2\beta'),\\
   \beta &= -\,\tfrac{1}{2 R}\,(8 S\alpha' + T\beta');
\end{aligned}
\]
so that the relation between $(\alpha, \beta)$ and $(\alpha', \beta')$ in the present case is similar to that between $(\alpha', \beta')$ and $(\alpha, \beta)$ in the former case. It may be noticed that in all these systems of linear equations, the determinant of transformation is a multiple of $64 S^3 - T^2\,(= R)$.

\newpage
\setcounter{page}{334}

\noindent\textbf{334}\hfill\textsc{a seventh memoir on quantics.}\hfill\textbf{[269}

\medskip

\noindent 245. \ It will be convenient, before giving the Tables for the covariants of
\[
   \alpha U + 6\beta H U,\quad 2\alpha C U - 2\beta D U,\quad 6\alpha P U + \beta Q U,\quad 2\alpha Y U - 2\beta Z U,
\]
which replace Tables 68 and 69 of my Third Memoir, to give the following separate Table of the quantics in $(\alpha, \beta)$ which enter into the expressions of the invariants in Tables 68 and 69, and in these new Tables.

\begin{center}Table No.~73.\end{center}
\[
\begin{aligned}
   &(1,\ 0,\ -24 S,\ -8 T,\ -48 S^2\Cdel\alpha,\,\beta)^4,\\
   &(S,\ T,\ 24 S^2,\ 4 T S,\ T^2 - 48 S^3\Cdel\alpha,\,\beta)^4,\\
   &(T,\ 96 S^2,\ 60 T S,\ 20 T^2,\ 240 T S^2,\ -48 T^2 S + 4608 S^4,\ -8 T^3 + 576 T^2 S^2\Cdel\alpha,\,\beta)^6.\\[6pt]
   &(48 S,\ 8 T,\ -96 S^2,\ -24 T S,\ -T^2 - 16 S^3\Cdel\alpha,\,\beta)^4,\\
   &(T^2 + 192 S^3,\ 128\,T S^2,\ 18\,T^2 S + 384\,S^4,\ 6 T^2,\ 6 T S\Cdel\alpha,\,\beta)^4,\\
   &(T^3 + 64\,T S^3,\ 5 T^2 S^2 - 64\,S^5,\ -\,8\,T^3 + 4608\,T S^4,\ 1920\,T^2 S^3 + 73728\,S^6,\\
   &\qquad 360\,T^3 S + 38400\,T S^4,\ 20\,T^4 + 8960\,T^2 S^3,\\
   &\qquad 840\,T^3 S^2 + 7680\,T S^5,\\
   &\qquad 36\,T^4 S + 384\,T^2 S^4 + 24576\,S^7,\\
   &\qquad T^5 - 40\,T^3 S^3 + 2560\,T S^6\Cdel\alpha,\,\beta)^6,
\end{aligned}
\]
where the first part of the Table contains the quantics in $(\alpha, \beta)$ which relate to the forms $\alpha U + 6\beta H U$ and $-\,2\alpha Y U + 2\beta Z U$, and the second part of the Table contains the quantics in $(\alpha, \beta)$ which relate to the forms $6\alpha P U + \beta Q U$ and $-\,2\alpha C U + 2\beta D U$.

\medskip

\noindent The quantics in $(\alpha, \beta)$ contained in the foregoing Table are in the sequel indicated by means of their leading coefficients; as thus,
\[
   (1,\,0,\,-24 S,\ldots\Cdel\alpha,\,\beta)^4,\quad (S,\,T,\ldots\Cdel\alpha,\,\beta)^4,\quad (T^2 + 192 S^3,\ldots\Cdel\alpha,\,\beta)^4,\ \&c.
\]

\medskip

\noindent 246. \ It is easy to see what transformations must be performed on the results in Tables 68 and 69, in order to obtain the new Tables. Thus, in the formation of Table 74, Table 68 gives $\alpha U + 6\beta H U$ and $H\,(\alpha U + 6\beta H U)$, and from these $C\,(\alpha U + 6\beta H U)$,

\newpage
\setcounter{page}{335}

\noindent\textbf{269]}\hfill\textsc{a seventh memoir on quantics.}\hfill\textbf{335}

\medskip

\noindent $D\,(\alpha U + 6\beta H U)$ have to be found: the same Table gives also $P\,(\alpha U + 6\beta H U)$, $Q\,(\alpha U + 6\beta H U)$, but the expressions of these quantities $Y U$, $Z U$ have to be introduced in the place of $P U$, $Q U$; and from the expressions so transformed are deduced also the expressions for $Y\,(\alpha U + 6\beta H U)$, $Z\,(\alpha U + 6\beta H U)$. Table 75 is to be deduced from Table 69 by writing therein $(\alpha, \beta')$, for $(\alpha, \beta)$, and then putting $6\alpha' P U + \beta' Q U = -\,2\alpha Y U + 2\beta Z U$ and $\alpha', \beta'$ but, as is seen above, if for $\alpha', \beta'$ as functions of $\alpha, \beta$ and of the invariants which is in some of the formul\ae, $Y U$, $Z U$, have to be introduced in the place of $P U$, $Q U$. And so for the Tables 76 and 77. The actual effectuation of the transformations would, it is almost needless to remark, be very laborious, but the forms of the results are easily foreseen, and the results can then be verified by means of one or two coefficients only. The new Tables are

\begin{center}Table No.~74.\end{center}
\[
\begin{aligned}
   R\,(\alpha U + 6\beta H U) &= R\times[(1,\,0,\,-24 S,\ldots\Cdel\alpha,\,\beta)^4]^3,\\
   S\,(\alpha U + 6\beta H U) &= (S,\,T,\ldots\Cdel\alpha,\,\beta)^4,\\
   T\,(\alpha U + 6\beta H U) &= [(T,\,96 S^2,\ldots\Cdel\alpha,\,\beta)^4]^2.\\[6pt]
   (\alpha U + 6\beta H U) &= \alpha U + 6\beta H U,\\
   H\,(\alpha U + 6\beta H U) &= -\,\tfrac{1}{2}\times\left\{\begin{aligned}&\partial_\beta\,(1,\,0,\,-24 S,\ldots\Cdel\alpha,\,\beta)^4\cdot U\\ &{} - 6\partial_\alpha\,(1,\,0,\,-24 S,\ldots\Cdel\alpha,\,\beta)^4\cdot H U,\end{aligned}\right.\\
   C\,(\alpha U + 6\beta H U) &= (1,\,0,\,-24 S,\ldots\Cdel\alpha,\,\beta)^4\,\times\\
                              &\quad\left\{\begin{aligned}&\partial_\beta\,(S,\,T,\ldots\Cdel\alpha,\,\beta)^4\cdot U\\ &{} - 6\partial_\alpha\,(S,\,T,\ldots\Cdel\alpha,\,\beta)^4\cdot H U,\end{aligned}\right.\\
   D\,(\alpha U + 6\beta H U) &= \tfrac{1}{12}\,(1,\,0,\,-24 S,\ldots\Cdel\alpha,\,\beta)^4\,\times\\
                              &\quad\left\{\begin{aligned}&\partial_\beta\,(T,\,96 S^2,\ldots\Cdel\alpha,\,\beta)^4\cdot U\\ &{} + \partial_\alpha\,(T,\,96 S^2,\ldots\Cdel\alpha,\,\beta)^4\cdot H U,\end{aligned}\right.\\
   P\,(\alpha U + 6\beta H U) &= -\,\tfrac{1}{3 R}\,\times\left\{\begin{aligned}&\partial_\beta\,(S,\,T,\ldots\Cdel\alpha,\,\beta)^4\cdot Y U\\ &{} + \partial_\alpha\,(S,\,T,\ldots\Cdel\alpha,\,\beta)^4\cdot Z U,\end{aligned}\right.\\
   Q\,(\alpha U + 6\beta H U) &= -\,\tfrac{2}{3 R}\,\times\left\{\begin{aligned}&\partial_\beta\,(T,\,96 S^2,\ldots\Cdel\alpha,\,\beta)^4\cdot Y U\\ &{} + \partial_\alpha\,(T,\,96 S^2,\ldots\Cdel\alpha,\,\beta)^4\cdot Z U,\end{aligned}\right.\\
   Y\,(\alpha U + 6\beta H U) &= [(1,\,0,\,-24 S,\ldots\Cdel\alpha,\,\beta)^4]^2\,\times\,(-\,2\alpha Y U + 2\beta Z U),\\
   Z\,(\alpha U + 6\beta H U) &= \tfrac{1}{2}\,[(1,\,0,\,-24 S,\ldots\Cdel\alpha,\,\beta)^4]^2\,\times\\
                              &\quad\left\{\begin{aligned}&\partial_\beta\,(1,\,0,\,-24 S,\ldots\Cdel\alpha,\,\beta)^4\cdot Y U\\ &{} + \partial_\alpha\,(1,\,0,\,-24 S,\ldots\Cdel\alpha,\,\beta)^4\cdot Z U.\end{aligned}\right.
\end{aligned}
\]

\newpage
\setcounter{page}{336}

\noindent\textbf{336}\hfill\textsc{a seventh memoir on quantics.}\hfill\textbf{[269}

\begin{center}No.~75.\end{center}
\[
\begin{aligned}
   R\,(-\,2\alpha Y U + 2\beta Z U) &= -\,4096 R^4\times[(S,\,T,\ldots\Cdel\alpha,\,\beta)^4]^3,\\
   S\,(-\,2\alpha Y U + 2\beta Z U) &= -\,R^2\times(1,\,0,\,-24 S,\ldots\Cdel\alpha,\,\beta)^4,\\
   T\,(-\,2\alpha Y U + 2\beta Z U) &= -\,8 R^3\times(T,\,96 S^2,\ldots\Cdel\alpha,\,\beta)^4.\\[6pt]
   -\,2\alpha Y U + 2\beta Z U &= -\,2\alpha Y U + 2\beta Z U,\\
   H\,(-\,2\alpha Y U + 2\beta Z U) &= -\,\tfrac{1}{2}\,R\,\times\left\{\begin{aligned}&\partial_\beta\,(S,\,T,\ldots\Cdel\alpha,\,\beta)^4\cdot Y U\\ &{} + \partial_\alpha\,(S,\,T,\ldots\Cdel\alpha,\,\beta)^4\cdot Z U,\end{aligned}\right.\\
   C\,(-\,2\alpha Y U + 2\beta Z U) &= -\,16 R^2\,(S,\,T,\ldots\Cdel\alpha,\,\beta)^4\,\times\\
                                    &\quad\left\{\begin{aligned}&\partial_\beta\,(1,\,0,\,-24 S,\ldots\Cdel\alpha,\,\beta)^4\cdot Y U\\ &{} + \partial_\alpha\,(1,\,0,\,-24 S,\ldots\Cdel\alpha,\,\beta)^4\cdot Z U,\end{aligned}\right.\\
   D\,(-\,2\alpha Y U + 2\beta Z U) &= -\,\tfrac{4}{3}\,R^2\,(S,\,T,\ldots\Cdel\alpha,\,\beta)^4\,\times\\
                                    &\quad\left\{\begin{aligned}&\partial_\beta\,(T,\,96 S^2,\ldots\Cdel\alpha,\,\beta)^4\cdot Y U\\ &{} + \partial_\alpha\,(T,\,96 S^2,\ldots\Cdel\alpha,\,\beta)^4\cdot Z U,\end{aligned}\right.\\
   P\,(-\,2\alpha Y U + 2\beta Z U) &= \tfrac{1}{3}\,R^2\,\times\left\{\begin{aligned}&\partial_\beta\,(1,\,0,\,-24 S,\ldots\Cdel\alpha,\,\beta)^4\cdot U\\ &{} - 6\partial_\alpha\,(1,\,0,\,-24 S,\ldots\Cdel\alpha,\,\beta)^4\cdot H U,\end{aligned}\right.\\
   Q\,(-\,2\alpha Y U + 2\beta Z U) &= -\,\tfrac{2}{3}\,R^2\,\times\left\{\begin{aligned}&\partial_\beta\,(T,\,96 S^2,\ldots\Cdel\alpha,\,\beta)^4\cdot U\\ &{} - 6\partial_\alpha\,(T,\,96 S^2,\ldots\Cdel\alpha,\,\beta)^4\cdot H U,\end{aligned}\right.\\
   Y\,(-\,2\alpha Y U + 2\beta Z U) &= 256 R^4\,[(S,\,T,\ldots\Cdel\alpha,\,\beta)^4]^2\,\times\,(\alpha U + 6\beta H U),\\
   Z\,(-\,2\alpha Y U + 2\beta Z U) &= -\,512 R^4\,[(S,\,T,\ldots\Cdel\alpha,\,\beta)^4]^2\,\times\\
                                    &\quad\left\{\begin{aligned}&\partial_\beta\,(S,\,T,\ldots\Cdel\alpha,\,\beta)^4\cdot U\\ &{} - 6\partial_\alpha\,(S,\,T,\ldots\Cdel\alpha,\,\beta)^4\cdot H U.\end{aligned}\right.
\end{aligned}
\]

\newpage
\setcounter{page}{337}

\noindent\textbf{269]}\hfill\textsc{a seventh memoir on quantics.}\hfill\textbf{337}

\begin{center}No.~76.\end{center}
\[
\begin{aligned}
   R\,(2\alpha C U - 2\beta D U) &= -\,4096 R^4\times[(T^2 + 192 S^3,\ldots\Cdel\alpha,\,\beta)^4]^3,\\
   S\,(2\alpha C U - 2\beta D U) &= -\,8 R^2\times(48 S,\,8 T,\ldots\Cdel\alpha,\,\beta)^4,\\
   T\,(2\alpha C U - 2\beta D U) &= -\,8 R^3\times(-\,8 T^3 + 4608 T S^4,\ldots\Cdel\alpha,\,\beta)^4.\\[6pt]
   2\alpha C U - 2\beta D U &= 2\alpha C U - 2\beta D U,\\
   H\,(2\alpha C U - 2\beta D U) &= \tfrac{1}{2}\times\left\{\begin{aligned}&\partial_\beta\,(T^2 + 192 S^3,\ldots\Cdel\alpha,\,\beta)^4\cdot C U\\ &{} + \partial_\alpha\,(T^2 + 192 S^3,\ldots\Cdel\alpha,\,\beta)^4\cdot D U,\end{aligned}\right.\\
   C\,(2\alpha C U - 2\beta D U) &= -\,16 R^2\,(T^2 + 192 S^3,\ldots\Cdel\alpha,\,\beta)^4\,\times\\
                                  &\quad\left\{\begin{aligned}&\partial_\beta\,(48 S,\,8 T,\ldots\Cdel\alpha,\,\beta)^4\cdot C U\\ &{} + \partial_\alpha\,(48 S,\,8 T,\ldots\Cdel\alpha,\,\beta)^4\cdot D U,\end{aligned}\right.\\
   D\,(2\alpha C U - 2\beta D U) &= \tfrac{4}{3}\,R^2\,(T^2 + 192 S^3,\ldots\Cdel\alpha,\,\beta)^4\,\times\\
                                  &\quad\left\{\begin{aligned}&\partial_\beta\,(-\,8 T^3 + 4608 T S^4,\ldots\Cdel\alpha,\,\beta)^4\cdot C U\\ &{} + \partial_\alpha\,(-\,8 T^3 + 4608 T S^4,\ldots\Cdel\alpha,\,\beta)^4\cdot D U,\end{aligned}\right.\\
   P\,(2\alpha C U - 2\beta D U) &= \tfrac{1}{6 R}\,\times\left\{\begin{aligned}&6\partial_\beta\,(48 S,\,8 T,\ldots\Cdel\alpha,\,\beta)^4\cdot P U\\ &{} + \partial_\alpha\,(48 S,\,8 T,\ldots\Cdel\alpha,\,\beta)^4\cdot Q U,\end{aligned}\right.\\
   Q\,(2\alpha C U - 2\beta D U) &= \tfrac{4}{R}\,\times\left\{\begin{aligned}&6\partial_\beta\,(-\,8 T^3 + 4608 T S^4,\ldots\Cdel\alpha,\,\beta)^4\cdot P U\\ &{} - \partial_\alpha\,(-\,8 T^3 + 4608 T S^4,\ldots\Cdel\alpha,\,\beta)^4\cdot Q U,\end{aligned}\right.\\
   Y\,(2\alpha C U - 2\beta D U) &= 216 R^2\,[(T^2 + 192 S^3,\ldots\Cdel\alpha,\,\beta)^4]^2\,\times\,(6\alpha P U + \beta Q U),\\
   Z\,(2\alpha C U - 2\beta D U) &= -\,512 R^4\,[(T^2 + 192 S^3,\ldots\Cdel\alpha,\,\beta)^4]^2\,\times\\
                                  &\quad\left\{\begin{aligned}&6\partial_\beta\,(T^2 + 192 S^3,\ldots\Cdel\alpha,\,\beta)^4\cdot P U\\ &{} - \partial_\alpha\,(T^2 + 192 S^3,\ldots\Cdel\alpha,\,\beta)^4\cdot Q U.\end{aligned}\right.
\end{aligned}
\]

\hfill C.~IV.\hfill 43

\newpage
\setcounter{page}{338}

\noindent\textbf{338}\hfill\textsc{a seventh memoir on quantics.}\hfill\textbf{[269}

\begin{center}No.~77.\end{center}
\[
\begin{aligned}
   R\,(6\alpha P U + \beta Q U) &= R^2\times[(48 S,\,8 T,\ldots\Cdel\alpha,\,\beta)^4]^3,\\
   S\,(6\alpha P U + \beta Q U) &= (T^2 + 192 S^3,\ldots\Cdel\alpha,\,\beta)^4,\\
   T\,(6\alpha P U + \beta Q U) &= (-\,8 T^3 + 4608 T S^4,\ldots\Cdel\alpha,\,\beta)^4.\\[6pt]
   (6\alpha P U + \beta Q U) &= 6\alpha P U + \beta Q U,\\
   H\,(6\alpha P U + \beta Q U) &= -\,\tfrac{1}{R}\,\times\left\{\begin{aligned}&-\,6\partial_\beta\,(48 S,\,8 T,\ldots\Cdel\alpha,\,\beta)^4\cdot P U\\ &{} - \partial_\alpha\,(48 S,\,8 T,\ldots\Cdel\alpha,\,\beta)^4\cdot Q U,\end{aligned}\right.\\
   C\,(6\alpha P U + \beta Q U) &= -\,(48 S,\,8 T,\ldots\Cdel\alpha,\,\beta)^4\,\times\\
                                 &\quad\left\{\begin{aligned}&-\,6\partial_\beta\,(T^2 + 192 S^3,\ldots\Cdel\alpha,\,\beta)^4\cdot P U\\ &{} - \partial_\alpha\,(T^2 + 192 S^3,\ldots\Cdel\alpha,\,\beta)^4\cdot Q U,\end{aligned}\right.\\
   D\,(6\alpha P U + \beta Q U) &= \tfrac{1}{12}\,(48 S,\,8 T,\ldots\Cdel\alpha,\,\beta)^4\,\times\\
                                 &\quad\left\{\begin{aligned}&\partial_\beta\,(-\,8 T^3 + 4608 T S^4,\ldots\Cdel\alpha,\,\beta)^4\cdot Q U,\end{aligned}\right.\\
   P\,(6\alpha P U + \beta Q U) &= -\,\tfrac{1}{R}\,\times\left\{\begin{aligned}&\partial_\beta\,(T^2 + 192 S^3,\ldots\Cdel\alpha,\,\beta)^4\cdot C U\\ &{} - \partial_\alpha\,(T^2 + 192 S^3,\ldots\Cdel\alpha,\,\beta)^4\cdot D U,\end{aligned}\right.\\
   Y\,(6\alpha P U + \beta Q U) &= -\,\tfrac{1}{2}\,R\,[(48 S,\,8 T,\ldots\Cdel\alpha,\,\beta)^4]^2\,\times\,(-\,2\alpha C U + 2\beta D U),\\
   Z\,(6\alpha P U + \beta Q U) &= -\,\tfrac{1}{4}\,R\,[(48 S,\,8 T,\ldots\Cdel\alpha,\,\beta)^4]^2\,\times\\
                                 &\quad\left\{\begin{aligned}&\partial_\beta\,(48 S,\,8 T,\ldots\Cdel\alpha,\,\beta)^4\cdot C U\\ &{} + \partial_\alpha\,(48 S,\,8 T,\ldots\Cdel\alpha,\,\beta)^4\cdot D U.\end{aligned}\right.
\end{aligned}
\]

\newpage
\setcounter{page}{339}

\noindent\textbf{269]}\hfill\textsc{a seventh memoir on quantics.}\hfill\textbf{339}

\medskip

\noindent 247. \ It will be noticed how Tables 74 and 75 form a system involving only the quantics in $(\alpha, \beta)$ contained in the first part of Table 73, and how, in like manner, Tables 76 and 77 form a system involving only the quantics in $(\alpha, \beta)$ contained in the second part of Table 73; and, moreover, how in each pair of Tables the covariants, \&c.\ correspond to each other as follows, viz.
\[
\begin{aligned}
   &R,\ S,\ T,\ 1,\ H,\ C,\ D,\ P,\ Q,\ Y,\ Z\quad\text{to}\\
   &S,\ R,\ T,\ Y,\ P,\ Z,\ Q,\ H,\ D,\ 1,\ C.
\end{aligned}
\]

\medskip

\noindent Thus in Table 74,---the formula for $H\,(\alpha U + 6\beta H U)$, and in Table 75,---the formula for $P\,(2\alpha Y U - 2\beta Z U)$, each of them involve the same factor
\[
   \left\{\begin{aligned}&\partial_\beta\,(1,\,0,\,-24 S,\ldots\Cdel\alpha,\,\beta)^4\cdot U\\ &{} - 6\partial_\alpha\,(1,\,0,\,-24 S,\ldots\Cdel\alpha,\,\beta)^4\cdot H U,\end{aligned}\right.
\]
and so in all the other cases.

\medskip

\noindent 248. \ The quantics in $(\alpha, \beta)$ in each part of the foregoing Table 73 are covariantively connected together. In fact, considering the function $(1,\,0,\,-24 S,\ldots\Cdel\alpha,\,\beta)^4$, which for shortness I call $G$, we have
\[
\begin{aligned}
   G &= (1,\,0,\,-24 S,\ldots\Cdel\alpha,\,\beta)^4,\\
   I G &= 0,\\
   J G &= 4\,(64 S^3 - T^2) = 4 R,\\
   D G &= (I G)^3 - 27\,(J G)^2 = -\,432 R^2,\\
   H G &= -\,\tfrac{1}{3}\,(S,\,T,\ldots\Cdel\alpha,\,\beta)^4,\\
   \Phi G &= 2\,(T,\,96 S^2,\ldots\Cdel\alpha,\,\beta)^6.
\end{aligned}
\]
The last-mentioned formul\ae, by the aid of Table 72, give rise to the following more general system in which they are themselves included.

\begin{center}Table No.~78.\end{center}
\[
\begin{aligned}
   \lambda G + 6\mu H G &= \lambda G + 6\mu H G,\\
   H\,(\lambda G + 6\mu H G) &= 36 R\mu^2 G + \lambda^2 H G,\\
   \Phi\,(\lambda G + 6\mu H G) &= (\lambda^3 - 216 R\mu^3)\cdot 2\,(T,\,96 S^2,\ldots\Cdel\alpha,\,\beta)^6,\\
   I\,(\lambda G + 6\mu H G) &= 72 R\mu^2,\\
   J\,(\lambda G + 6\mu H G) &= 4 R\,(\lambda^3 + 216 R\mu^3),\\
   D\,(\lambda G + 6\mu H G) &= -\,432 R^2\,(\lambda^3 - 216 R\mu^3)^2.
\end{aligned}
\]

\hfill 43--2

\newpage
\setcounter{page}{340}

\noindent\textbf{340}\hfill\textsc{a seventh memoir on quantics.}\hfill\textbf{[269}

\medskip

\noindent The expression for $H\,(\lambda G + 6\mu H G)$, putting therein $\lambda = 0$, shows that, to a numerical factor \textit{pr\`es}, $H\cdot H G$ is equal to $G$, and hence, disregarding numerical factors, we may say that each of the quartics $(1,\,0,\,-24 S,\ldots\Cdel\alpha,\,\beta)^4$, $(S,\,T,\ldots\Cdel\alpha,\,\beta)^4$ is the Hessian of the other of them, and that the sextic $(T,\,96 S^2,\ldots\Cdel\alpha,\,\beta)^6$ is the cubicovariant of each of them.

\medskip

\noindent 249. \ Similarly, if the function $(48 S,\,8 T,\ldots\Cdel\alpha,\,\beta)^4$ is for shortness called $G$, then we have
\[
\begin{aligned}
   I G &= 0,\\
   J G &= 4\,(64 S^3 - T^2)^2 = 4 R,\\
   D G &= (I G)^3 - 27\,(J G)^2 = -\,432 R^2,\\
   H G &= -\,\tfrac{1}{3}\,(T^2 + 192 S^3,\ldots\Cdel\alpha,\,\beta)^4,\\
   \Phi G &= -\,2\,(-\,8 T^3 + 4608 T S^4,\ldots\Cdel\alpha,\,\beta)^6.
\end{aligned}
\]
The last-mentioned formul\ae, by the aid of the same Table 72, give rise to the more general system in which they are themselves included.

\begin{center}Table No.~79.\end{center}
\[
\begin{aligned}
   \lambda G + 6\mu H G &= \lambda G + 6\mu H G,\\
   H\,(\lambda G + 6\mu H G) &= 36 R^2\mu^2 G + \lambda^2 H G,\\
   \Phi\,(\lambda G + 6\mu H G) &= (\lambda^3 - 216 R^2\mu^3)\times -\,2\,(-\,8 T^3 + 4608 T S^4,\ldots\Cdel\alpha,\,\beta)^6,\\
   I\,(\lambda G + 6\mu H G) &= 72 R^2\mu^2,\\
   J\,(\lambda G + 6\mu H G) &= 4 R\,(\lambda^3 + 216 R^2\mu^3),\\
   D\,(\lambda G + 6\mu H G) &= -\,432 R^2\,(\lambda^3 - 216 R^2\mu^3)^2.
\end{aligned}
\]

\medskip

\noindent The expression for $H\,(\lambda G + 6\mu H G)$, putting therein $\lambda = 0$, shows that, to a numerical factor \textit{pr\`es}, $H\cdot H G$ is equal to $G$; so that, disregarding numerical factors, we may say that each of the quartics $(48 S,\,T,\ldots\Cdel\alpha,\,\beta)^4$, $(T^2 + 192 S^3,\ldots\Cdel\alpha,\,\beta)^4$ is the Hessian of the other of them, and that the sextic $(-\,8 T^3 + 4608 T S^4,\ldots\Cdel\alpha,\,\beta)^6$ is the cubicovariant of each of them.

\medskip

\noindent 250. \ But besides this, the quantics in $(\alpha, \beta)$ in the two parts of the Table 73 are linearly connected together: the linear relations in question are in fact the equations whereon depend the expressions for the invariants in Tables 76 and 77 as deduced from those in Tables 74 and 75; and in the order of proof, they precede the formul\ae{} in these four Tables. The linear relations are

\newpage
\setcounter{page}{341}

\noindent\textbf{269]}\hfill\textsc{a seventh memoir on quantics.}\hfill\textbf{341}

\begin{center}Table No.~80.\end{center}
\[
\begin{aligned}
   (1,\ 0,\ -24 S,\ldots\Cdel -\,2 T\alpha - 16 S^2\beta,\ 8 S\alpha + T\beta)^4 &= -\,16 R\,(T^2 + 192 S^3,\ldots\Cdel\alpha,\,\beta)^4,\\
   (S,\ T,\ldots\Cdel -\,2 T\alpha - 16 S^2\beta,\ 8 S\alpha + T\beta)^4 &= R^2\,(48 S,\,8 T,\ldots\Cdel\alpha,\,\beta)^4,\\
   (T,\ 96 S^2,\ldots\Cdel -\,2 T\alpha - 16 S^2\beta,\ 8 S\alpha + T\beta)^6 &= -\,8 R^2\,(-\,8 T^3 + 4608 T S^4,\ldots\Cdel\alpha,\,\beta)^6.\\[6pt]
   (48 S,\ 8 T,\ldots\Cdel T\alpha + 16 S^2\beta,\ -\,8 S\alpha - 2 T\beta)^4 &= -\,16 R^2\,(S,\,T,\ldots\Cdel\alpha,\,\beta)^4,\\
   (T^2 + 192 S^3,\ldots\Cdel T\alpha + 16 S^2\beta,\ -\,8 S\alpha - 2 T\beta)^4 &= -\,R^2\,(1,\,0,\,-24 S,\ldots\Cdel\alpha,\,\beta)^4,\\
   (-\,8 T^3 + 4608 T S^4,\ldots\Cdel T\alpha + 16 S^2\beta,\ -\,8 S\alpha - 2 T\beta)^6 &= -\,8 R^3\,(T,\,96 S^2,\ldots\Cdel\alpha,\,\beta)^6.
\end{aligned}
\]

\medskip

\noindent Hence, attending to the remarks on the Tables 78 and 79, we may say that the quartics
\[
   (1,\,0,\,-24 S,\ldots\Cdel\alpha,\,\beta)^4,\quad (48 S,\,T,\ldots\Cdel\alpha,\,\beta)^4,
\]
which belong to the two parts respectively of Table 73, and which are related, the first of them to the discriminants of $\alpha U + 6\beta H U$ and $-\,2\alpha Y U + 2\beta Z U$, and the second to the discriminants of $6\alpha P U + \beta Q U$, $-\,2\alpha C U + 2\beta D U$, have these relations to each other, viz.\ each is a linear transformation of the Hessian of the other of them, and the cubicovariant of each is a linear transformation of the cubicovariant of the other.

\newpage
\setcounter{page}{342}

% ============================================================
% PAPER 270
% ON THE DOUBLE TANGENTS OF A CURVE OF THE FOURTH ORDER
% ============================================================

\begin{center}
\textbf{270.}\\[1em]
{\Large ON THE DOUBLE TANGENTS OF A CURVE OF THE FOURTH}\\[0.3em]
{\Large ORDER.}\\[1em]
\textit{[From the Philosophical Transactions of the Royal Society of London, vol.~\textsc{cli}.\ (for the year 1861), pp.~357--362.\ \ Received May 30,---Read June 20, 1861.]}
\end{center}

\medskip

\noindent\textsc{The} present memoir is intended to be supplementary to that ``On the Double Tangents of a Plane Curve'' (\textit{Phil.\ Trans.}, vol.~\textsc{cxlix}.\ (1859), pp.~193--212) [260]. I take the opportunity of correcting an error which I have there fallen into, and which is rather a misleading one, viz.\ the emanants $U_1, U_2,\ldots$ were numerically determined in such manner as to become equal to $U$ on putting $(x_1, y_1, z_1)$ equal to $(x, y, z)$; the numerical determination should have been (and in the latter part of the memoir is assumed to be) such as to render $H_1, H_2$, \&c.\ equal to $H$, on making the substitution in question; that is, in the place of the formul\ae
\[
\begin{aligned}
   U_1 &= \tfrac{1}{n}\,(x_1\partial_x + y_1\partial_y + z_1\partial_z)\,U,\\
   U_2 &= \tfrac{1}{n\,(n-1)}\,(x_1\partial_x + y_1\partial_y + z_1\partial_z)^2\,U,\ \&c.,
\end{aligned}
\]
there ought to have been
\[
\begin{aligned}
   U_1 &= \tfrac{1}{(n-2)}\,(x_1\partial_x + y_1\partial_y + z_1\partial_z)\,U,\\
   U_2 &= \tfrac{1}{(n-2)(n-3)}\,(x_1\partial_x + y_1\partial_y + z_1\partial_z)^2\,U,\ \&c.
\end{aligned}
\]
[this error is corrected \textit{ante} p.~189].

\medskip

\noindent The points of contact of the double tangents of the curve of the fourth order or quartic $U = 0$, are given as the intersections of the curve with a curve of the fourteenth order $\Pi = 0$; the last-mentioned curve is not absolutely determinate, since instead of $\Pi = 0$, we may, it is clear, write $\Pi + M U = 0$, where $M$ is an arbitrary

\newpage
\setcounter{page}{343}

\noindent\textbf{270]}\hfill\textsc{on the double tangents of a curve of the fourth order.}\hfill\textbf{343}

\medskip

\noindent function of the tenth order. I have in the memoir spoken of Hesse's original form (say $\Pi_1 = 0$) of the curve of the fourteenth order obtained by him in 1850, and of his transformed form (say $\Pi_2 = 0$) obtained in 1856. The method in the memoir itself (Mr~Salmon's method) gives, in the case in question of a quartic curve, a third form, say $\Pi_3 = 0$. It appears by his paper ``On the Determination of the Points of Contact of Double Tangents to an Algebraic Curve'' (\textit{Quart.\ Math.\ Journ.}\ vol.~\textsc{iii}.\ p.~317 (1859)), that Mr~Salmon has verified by algebraic transformations the equivalence of the last-mentioned form with those of Hesse; but the process is not given. The object of the present memoir is to demonstrate the equivalence in question, viz.\ that of the equation $\Pi_3 = 0$ with the one or other of the equations $\Pi_1 = 0$, $\Pi_2 = 0$, in virtue of the equation $U = 0$. The transformation depends, 1st, on a theorem used by Hesse for the deduction of his second form $\Pi_2 = 0$ from the original form $\Pi_1 = 0$, which theorem is given in his paper ``Transformation der Gleichung der Curven 14ten Grades welche eine gegebene Curve 4ten Grades in den Ber\"uhrungspuncten ihrer Doppeltangenten schneiden,'' \textit{Crelle}, t.~\textsc{lii}.\ pp.~97--103 (1856), containing the transformation in question; I prove this theorem in a different and (as it appears to me) more simple manner; 2nd, on a theorem relating to a cubic curve proved incidentally in my memoir ``On the Conic of Five-pointic Contact at any point of a Plane Curve'' (\textit{Phil.\ Trans.}, vol.~\textsc{cxlix}.\ (1859), see p.~385 [261]), the cubic curve being in the present case any first emanent of the given quartic curve: the demonstration occupies only a single paragraph, and it is here reproduced; and I reproduce also Hesse's demonstration of the equivalence of the two forms $\Pi_1 = 0$ and $\Pi_2 = 0$.

\medskip

\noindent Let $U = (*\Cdel x, y, z)^4$ be a quartic function of $(x, y, z)$; $(A, B, C, F, G, H)$ the reciprocal system
\[
   (b c - f^2,\ c a - g^2,\ a b - h^2,\ g h - a f,\ h f - b g,\ f g - c h);
\]
and let $H$ be the Hessian of $U$, or determinant $a b c - a f^2 - b g^2 - c h^2 + 2 f g h$ ($H$ is of course a sextic function of $x, y, z$); $(a', b', c', f', g', h')$ the second differential coefficients of $H$; $(A', B', C', F', G', H')$ the reciprocal system
\[
   (b' c' - f'^2,\ c' a' - g'^2,\ a' b' - h'^2,\ g' h' - a' f',\ h' f' - b' g',\ f' g' - c' h').
\]
Then $U = 0$ being the equation of a quartic curve, the equation of the curve of the fourteenth order which by its intersections determines the points of contact of the double tangents of the quartic curve, may be taken to be (Hesse's original form)
\[
   \Pi_1 = (A, B, C, F, G, H\Cdel\partial_x H,\partial_y H,\partial_z H)^2 - 3 H\,(A, B, C, F, G, H\Cdel\partial_x,\partial_y,\partial_z)^2\,H = 0.\footnote{In quoting this formula in my former memoir, the numerical factor 3 is by mistake omitted. [This correction should have been made \textit{ante} p.~187.]}
\]
Or it may be taken to be (Hesse's transformed form)
\[
   \Pi_2 = 5\,(A, B, C, F, G, H\Cdel\partial_x H,\partial_y H,\partial_z H)^2 - 3\,(A', B', C', F', G', H'\Cdel\partial_x U,\partial_y U,\partial_z U)^2 = 0.
\]
And moreover, if $U_1 = \tfrac{1}{2}\,(x_1\partial_x + y_1\partial_y + z_1\partial_z)\,U$, and if $H_1$ be the Hessian of $U_1$, and $(a'', b'', c'', f'', g'', h'')$ the second differential coefficients of $H - 3 H_1$, where in the differentiations $(x_1, y_1, z_1)$ are treated as constants but after the differentiations are

\newpage
\setcounter{page}{344}

\noindent\textbf{344}\hfill\textsc{on the double tangents of a curve of the fourth order.}\hfill\textbf{[270}

\medskip

\noindent effected they are replaced by $(x, y, z)$, and if $(A'', B'', C'', F'', G'', H'')$ be the reciprocal system
\[
   (b'' c'' - f''^2,\ c'' a'' - g''^2,\ a'' b'' - h''^2,\ g'' h'' - a'' f'',\ h'' f'' - b'' g'',\ f'' g'' - c'' h''),
\]
then the equations of the curve of the fourteenth order may be taken to be (Salmon's form)
\[
   \Pi_3 = (A'', B'', C'', F'', G'', H''\Cdel\partial_x U,\partial_y U,\partial_z U)^2 = 0.
\]

\medskip

\noindent I have preferred to write the three equations in the foregoing forms; but it is clear that the terms
\[
   (A, B, C, F, G, H\Cdel\partial_x,\partial_y,\partial_z)^2\,H\,;\quad (A', B', C', F', G', H'\Cdel\partial_x,\partial_y,\partial_z)^2\,U
\]
might also have been written
\[
   (A, B, C, F, G, H\Cdel a', b', c', 2 f', 2 g', 2 h');\quad (A', B', C', F', G', H'\Cdel a, b, c, 2 f, 2 g, 2 h).
\]

\medskip

\noindent As already noticed, it has been shown by Hesse (and his demonstration is to be here reproduced) that the two forms $\Pi_1 = 0$ and $\Pi_2 = 0$ are equivalent to each other. And the object of the memoir is to show that the third form $\Pi_3 = 0$ is equivalent to the other two. The equivalences in question subsist in virtue of the equation $U = 0$, that is, the functions $\Pi_1, \Pi_2, \Pi_3$ are not identical, but differ from each other by multiples of $U$.

\medskip

\begin{center}\textit{Demonstration of Hesse's Theorem.}\end{center}

\medskip

\noindent Let $(a, b, c, f, g, h)$, $(a', b', c', f', g', h')$ be any systems of coefficients of a ternary quadratic function; $(A, B, C, F, G, H)$, $(A', B', C', F', G', H')$ the reciprocal systems as above, $(x, y, z)$ arbitrary quantities. Consider the function
\[
\begin{aligned}
   \Omega &= (a, b, c, f, g, h\Cdel x, y, z)^2\,(A', B', C', F', G', H'\Cdel a, b, c, 2 f, 2 g, 2 h)\\
   &\quad{} - (A', B', C', F', G', H'\Cdel a x + h y + g z,\ h x + b y + f z,\ g x + f y + c z)^2.
\end{aligned}
\]

\medskip

\noindent The term involving $A'$ is
\[
   a\,(a, b, c, f, g, h\Cdel x, y, z)^2 - (a x + h y + g z)^2,
\]
which is
\[
\begin{aligned}
   &= (a b - h^2)\,y^2 + (a c - g^2)\,z^2 + 2\,(a f - g h)\,y z,\\
   &= C y^2 + B z^2 - 2 F y z;
\end{aligned}
\]
and the term involving $2 F'$ is
\[
   f\,(a, b, c, f, g, h\Cdel x, y, z)^2 - (h x + b y + f z)(g x + f y + c z),
\]
which is
\[
\begin{aligned}
   &= (a f - g h)\,x^2 + (f^2 - b c)\,y z + (f g - c h)\,z x + (h f - b g)\,x y,\\
   &= -\,F x^2 - A y z + H z x + G x y;
\end{aligned}
\]

\newpage
\setcounter{page}{345}

\noindent\textbf{270]}\hfill\textsc{on the double tangents of a curve of the fourth order.}\hfill\textbf{345}

\medskip

\noindent and the entire expression for $\Omega$ is thus
\[
\begin{aligned}
   &A'\,(C y^2 + B z^2 - 2 F y z)\\
   &{} + B'\,(A z^2 + C x^2 - 2 G z x)\\
   &{} + C'\,(B x^2 + A y^2 - 2 H x y)\\
   &{} + 2 F'\,(-\,F x^2 - A y z + H z x + G x y)\\
   &{} + 2 G'\,(-\,G y^2 - B z x + F x y + H y z)\\
   &{} + 2 H'\,(-\,H z^2 - C x y + G y z + F z x);
\end{aligned}
\]
or, what is the same thing,
\[
   \Omega = (B C' + B' C - 2 F F',\ C A' + C' A - 2 G G',\ A B' + A' B - 2 H H',\ G H' + G' H - A F' - A' F,\ H F' + H' F - B G' - B' G,\ F G' + F' G - C H' - C' H\Cdel x, y, z)^2,
\]
which is really the fundamental theorem. It is however used as follows; viz.\ the right-hand side being symmetrical in regard to the two systems
\[
   (a, b, c, f, g, h),\quad (a', b', c', f', g', h'),
\]
the left-hand side, which is not in form symmetrical as regards the two systems, must be so in reality; or if $\Omega'$ is what $\Omega$ becomes by interchanging the two systems, then $\Omega' = \Omega$; or substituting for $\Omega$ and $\Omega'$ their values, we have
\[
\begin{aligned}
   &(a, b, c, f, g, h\Cdel x, y, z)^2\cdot(A', B', C', F', G', H'\Cdel a, b, c, 2 f, 2 g, 2 h)\\
   &\quad{} - (A', B', C', F', G', H'\Cdel a x + h y + g z,\ h x + b y + f z,\ g x + f y + c z)^2\\
   &= (a', b', c', f', g', h'\Cdel x, y, z)^2\cdot(A, B, C, F, G, H\Cdel a', b', c', 2 f', 2 g', 2 h')\\
   &\quad{} - (A, B, C, F, G, H\Cdel a' x + h' y + g' z,\ h' x + b' y + f' z,\ g' x + f' y + c' z)^2,
\end{aligned}
\]
which is Hesse's theorem.

\medskip

\noindent If in particular $(a, b, c, f, g, h)$ are the second differential coefficients of a function $u = (*\Cdel x, y, z)^p$, and $(a', b', c', f', g', h')$ the second differential coefficients of a function $u' = (*\Cdel x, y, z)^{p'}$, then the equation becomes
\[
\begin{aligned}
   &p\,(p - 1)\,u\cdot(A', B', C', F', G', H'\Cdel\partial_x,\partial_y,\partial_z)^2\,u - (p - 1)^2\,(A', B', C', F', G', H'\Cdel\partial_x u,\partial_y u,\partial_z u)^2\\
   &= p'\,(p' - 1)\,u'\cdot(A, B, C, F, G, H\Cdel\partial_x,\partial_y,\partial_z)^2\,u' - (p' - 1)^2\,(A, B, C, F, G, H\Cdel\partial_x u',\partial_y u',\partial_z u')^2;
\end{aligned}
\]
and if for $u, u'$ we take the quartic function $U$ and the sextic function $H$, its Hessian, we have
\[
\begin{aligned}
   &12\,U\cdot(A', B', C', F', G', H'\Cdel\partial_x,\partial_y,\partial_z)^2\,U - 9\,(A', B', C', F', G', H'\Cdel\partial_x U,\partial_y U,\partial_z U)^2\\
   &= 30\,H\cdot(A, B, C, F, G, H\Cdel\partial_x,\partial_y,\partial_z)^2\,H - 25\,(A, B, C, F, G, H\Cdel\partial_x H,\partial_y H,\partial_z H)^2;
\end{aligned}
\]
and if in this identical equation we write $U = 0$, then from the resulting equation and the equation
\[
   \Pi_1 = -\,3 H\,(A, B, C, F, G, H\Cdel\partial_x,\partial_y,\partial_z)^2\,H + (A, B, C, F, G, H\Cdel\partial_x H,\partial_y H,\partial_z H)^2
\]
\hfill 44

\newpage
\setcounter{page}{346}

\noindent\textbf{346}\hfill\textsc{on the double tangents of a curve of the fourth order.}\hfill\textbf{[270}

\medskip

\noindent we may eliminate any one of the three terms
\[
\begin{aligned}
   &(A', B', C', F', G', H'\Cdel\partial_x U,\partial_y U,\partial_z U)^2,\\
   &H\cdot(A, B, C, F, G, H\Cdel\partial_x,\partial_y,\partial_z)^2\,H,\\
   &(A, B, C, F, G, H\Cdel\partial_x H,\partial_y H,\partial_z H)^2;
\end{aligned}
\]
and in particular if the second term be eliminated, we obtain the equation
\[
   \Pi_2 = 5\,(A, B, C, F, G, H\Cdel\partial_x H,\partial_y H,\partial_z H)^2 - 3\,(A', B', C', F', G', H'\Cdel\partial_x U,\partial_y U,\partial_z U)^2,
\]
and the equivalence of the two forms $\Pi_1 = 0$ and $\Pi_2 = 0$ is thus established.

\medskip

\noindent But Hesse's theorem leads also to the demonstration of the equivalence of the third form $\Pi_3 = 0$. To use it for this purpose, I remark that if $(a'', b'', c'', f'', g'', h'')$ are the second differential coefficients of $H - 3 H_1$, where after the differentiations $(x_1, y_1, z_1)$ are to be replaced by $(x, y, z)$, then the theorem gives
\[
\begin{aligned}
   &12\,U\cdot(A'', B'', C'', F'', G'', H''\Cdel\partial_x,\partial_y,\partial_z)^2\,U - 9\,(A'', B'', C'', F'', G'', H''\Cdel\partial_x U,\partial_y U,\partial_z U)^2\\
   &= (a'', b'', c'', f'', g'', h''\Cdel x, y, z)^2\cdot(A, B, C, F, G, H\Cdel\partial_x,\partial_y,\partial_z)^2\,(H - 3 H_1)\\
   &\quad{} - (A, B, C, F, G, H\Cdel a'' x + h'' y + g'' z,\ h'' x + b'' y + f'' z,\ g'' x + f'' y + c'' z)^2.
\end{aligned}
\]
But on putting $(x, y, z)$ for $(x_1, y_1, z_1)$ we have (since $H$ is a homogeneous function of the order 6, and $H_1$ before the change is a homogeneous function of the order 3 in $(x, y, z)$): $a'' x + h'' y + g'' z = 5\partial_x H - 3\cdot 2\,\partial_x H_1 = 5\partial_x H - 3\partial_x H$ (since, on making the substitution, $H_1 = H$, but $\partial_x H_1 = \tfrac{1}{2}\partial_x H$) $= 2\partial_x H$; and thus
\[
   (a'' x + h'' y + g'' z,\ h'' x + b'' y + f'' z,\ g'' x + f'' y + c'' z) = (2\partial_x H, 2\partial_y H, 2\partial_z H);
\]
and similarly, on making the substitution,
\[
   (a'', b'', c'', f'', g'', h''\Cdel x, y, z)^2 = 6\cdot 5 H - 3\cdot 3\cdot 2 H_1 = (30 - 18)\,H = 12\,H.
\]
Hence writing therein $U = 0$, the foregoing equation becomes
\[
\begin{aligned}
   -\,9\,(A'', B'', C'', F'', G'', H''\Cdel\partial_x U,\partial_y U,\partial_z U)^2 &= 12\,H\cdot(A, B, C, F, G, H\Cdel\partial_x,\partial_y,\partial_z)^2\,(H - 3 H_1)\\
   &\quad{} - 4\,(A, B, C, F, G, H\Cdel\partial_x H,\partial_y H,\partial_z H)^2,
\end{aligned}
\]
which may also be written
\[
\begin{aligned}
   -\,9\,(A'', B'', C'', F'', G'', H''\Cdel\partial_x U,\partial_y U,\partial_z U)^2 &= 12\,H\cdot(A, B, C, F, G, H\Cdel\partial_x,\partial_y,\partial_z)^2\,H\\
   &\quad{} - 36\,H\cdot(A, B, C, F, G, H\Cdel\partial_x,\partial_y,\partial_z)^2\,H_1\\
   &\quad{} - 4\,(A, B, C, F, G, H\Cdel\partial_x H,\partial_y H,\partial_z H)^2,
\end{aligned}
\]
where $(x_1, y_1, z_1)$ are ultimately to be replaced by $(x, y, z)$. The second line in fact vanishes, which I show as follows:

\newpage
\setcounter{page}{347}

\noindent\textbf{270]}\hfill\textsc{on the double tangents of a curve of the fourth order.}\hfill\textbf{347}

\begin{center}\textit{Demonstration of my Theorem for a Cubic Curve.}\end{center}

\medskip

\noindent Let $U = (*\Cdel x, y, z)^3$ be a cubic function; it may by a linear transformation of the coordinates be reduced to the canonical form $x^3 + y^3 + z^3 + 6 l x y z$, and we then have
\[
\begin{aligned}
(A, B, C, F, G, H\Cdel\partial_x, \partial_y, \partial_z)^2\,H \div 6^2\quad &= \quad(y z - l^2 x^2)\cdot -\,6 l^2 x\\
&\quad{} + (z x - l^2 y^2)\cdot -\,6 l^2 y\\
&\quad{} + (x y - l^2 z^2)\cdot -\,6 l^2 z\\
&\quad{} + 2\,(l^2 y z - l x^2)\cdot(1 + 2 l^3)\,x\\
&\quad{} + 2\,(l^2 z x - l y^2)\cdot(1 + 2 l^3)\,y\\
&\quad{} + 2\,(l^2 x y - l z^2)\cdot(1 + 2 l^3)\,z\\
&= -\,18 l^2 x y z + 6 l^4\,(x^3 + y^3 + z^3)\\
&\quad{} + 6 l^2\,(1 + 2 l^3)\,x y z - 2 l\,(1 + 2 l^3)\,(x^3 + y^3 + z^3)\\
&= (-\,12 l^2 + 12 l^5)\,x y z + (-\,2 l + 2 l^4)\,(x^3 + y^3 + z^3)\\
&= 2\,(-\,l + l^4)\,(x^3 + y^3 + z^3 + 6 l x y z);
\end{aligned}
\]
or since $-\,l + l^4$ is equal to the quartinvariant $S$, and the equation is an invariantive one, we have for any cubic function whatever
\[
   (A, B, C, F, G, H\Cdel\partial_x, \partial_y, \partial_z)^2\,H \div 6^2 = 2 S\cdot U,
\]
which is the theorem in question. There is a difference of notation, and consequently a different numerical factor, in the theorem as stated in the memoir on the conic of five-pointic contact, referred to above.

\medskip

\noindent If $U$ is as above, $U$ is a quartic function $(*\Cdel x, y, z)^4$, and $U_1 = \tfrac{1}{2}(x_1\partial_x + y_1\partial_y + z_1\partial_z)\,U$ is a cubic function, and we have
\[
   (A_1, B_1, C_1, F_1, G_1, H_1\Cdel\partial_x, \partial_y, \partial_z)^2\,H_1 \div 6^2 = 2 S_1\cdot U_1,
\]
where it is to be noticed that $S_1$ denotes a quartic function in the coefficients of $U_1$, and consequently a quartic function in $(x_1, y_1, z_1)$, the coefficients being quartic functions of the coefficients of $U$. On writing $(x, y, z)$ in the place of $(x_1, y_1, z_1)$, $S_1$ becomes a quartic function of $(x, y, z)$, which is in fact a quarticovariant quartic of $U$.

\medskip

\noindent If in the foregoing equation we write $(x, y, z)$ in the place of $(x_1, y_1, z_1)$, then $U_1$ becomes equal to $U$, and consequently, if $U = 0$, the right-hand side of the equation vanishes. Moreover, $(a_1, b_1, c_1, f_1, g_1, h_1)$, the second differential coefficients of $U_1$, become

\hfill 44--2

\newpage
\setcounter{page}{348}

\noindent\textbf{348}\hfill\textsc{on the double tangents of a curve of the fourth order.}\hfill\textbf{[270}

\medskip

\noindent equal to $(a, b, c, f, g, h)$, and consequently the coefficients $(A_1, B_1, C_1, F_1, G_1, H_1)$ become equal to $(A, B, C, F, G, H)$. Hence, assuming always that $U = 0$, the equation becomes
\[
   (A, B, C, F, G, H\Cdel\partial_x, \partial_y, \partial_z)^2\,H_1 = 0,
\]
where after the differentiations $(x_1, y_1, z_1)$ are replaced by $(x, y, z)$. This is the form which is required for the present purpose.

\medskip

\noindent Now returning to the foregoing expression of $-\,9\,(A'', B'', C'', F'', G'', H''\Cdel\partial_x U,\partial_y U,\partial_z U)^2$, this now becomes
\[
\begin{aligned}
   -\,9\,\Pi_3 &= -\,9\,(A'', B'', C'', F'', G'', H''\Cdel\partial_x U,\partial_y U,\partial_z U)^2\\
                &= 4\,\{3\,H\cdot(A, B, C, F, G, H\Cdel\partial_x,\partial_y,\partial_z)^2\,H - (A, B, C, F, G, H\Cdel\partial_x U,\partial_y U,\partial_z U)^2\},
\end{aligned}
\]
so that the equation $\Pi_3 = 0$ gives
\[
   \Pi_1 = (A, B, C, F, G, H\Cdel\partial_x U,\partial_y U,\partial_z U)^2 - 3 H\cdot(A, B, C, F, G, H\Cdel\partial_x,\partial_y,\partial_z)^2\,H = 0,
\]
and the equivalence of the equations $\Pi_1 = 0$ and $\Pi_3 = 0$ is thus established.

\newpage
\setcounter{page}{349}

% ============================================================
% PAPER 271
% SUR L'INVARIANT LE PLUS SIMPLE D'UNE FONCTION QUAD-
% RATIQUE BI-TERNAIRE
% ============================================================

\selectlanguage{french}

\begin{center}
\textbf{271.}\\[1em]
{\Large SUR L'INVARIANT LE PLUS SIMPLE D'UNE FONCTION QUAD-}\\[0.3em]
{\Large RATIQUE BI-TERNAIRE, ET SUR LE R\'ESULTANT DE TROIS}\\[0.3em]
{\Large FONCTIONS QUADRATIQUES TERNAIRES.}\\[1em]
\textit{[From the Journal f\"ur die reine und angewandte Mathematik (Crelle), tom.~\textsc{lvii}.\ (1860), pp.~139--148.]}
\end{center}

\medskip

\noindent\textsc{M.~Sylvester} a trouv\'e depuis longtemps, pour le r\'esultant des trois fonctions quadratiques ternaires
\[
\begin{aligned}
   &(a,\ b,\ c,\ f,\ g,\ h)\,(x, y, z)^2,\\
   &(a',\ b',\ c',\ f',\ g',\ h')\,(x, y, z)^2,\\
   &(a'',\ b'',\ c'',\ f'',\ g'',\ h'')\,(x, y, z)^2,
\end{aligned}
\]
une expression remarquable
\[
   12 R = 16 C_{12} - C_6{}^2,
\]
o\`u $C_{12}, C_6$ repr\'esentent des fonctions donn\'ees des coefficients qui sont non seulement des invariants mais aussi des combinants des trois fonctions quadratiques (Voir le m\'emoire de M.~Sylvester, ``On the Calculus of forms otherwise the theory of Invariants,'' \S~VII., \textit{Camb.\ et Dub.\ Math.\ Journ.}, t.~\textsc{viii}., de la, pp.~256--269 de l'ann\'ee 1853, \'Equation (A.), p.~267).

\medskip

\noindent Pour expliquer la formation du fonction $C_6$, il convient de consid\'erer la fonction quadratique bi-ternaire form\'ee par les coefficients $a, b, c, f, g, h$; $a', b', c', f', g', h'$; $a'', b'', c'', f'', g'', h''$ et par les coefficients du quadratique r\'eciproque
\[
   A,\ B,\ C,\ F,\ G,\ H,\quad A',\ B',\ C',\ F',\ G',\ H',\quad A'',\ B'',\ C'',\ F'',\ G'',\ H''\ ;
\]
savoir je consid\`ere
\[
   U = \left|\begin{array}{cccccc} a, & h, & g, & F, & G, & H\\ h, & b, & f, & F', & G', & H'\\ g, & f, & c, & F'', & G'', & H''\\ F, & F', & F'', & a', & h', & g'\\ G, & G', & G'', & h', & b', & f'\\ H, & H', & H'', & g', & f', & c'\end{array}\right|\cdot(x, y, z, \xi, \eta, \zeta)^2
\]

\newpage
\setcounter{page}{350}

\noindent\textbf{350}\hfill\textsc{sur l'invariant le plus simple d'une fonction \&c.}\hfill\textbf{[271}

\medskip

\noindent cette notation repr\'esentant la fonction
\[
\begin{aligned}
   &(a,\ b,\ c,\ f,\ g,\ h)\,(x, y, z)^2\cdot\xi^2\\
   &\quad{} + (a',\ b',\ c',\ f',\ g',\ h')\,(x, y, z)^2\cdot\eta^2\\
   &\quad{} + (a'',\ b'',\ c'',\ f'',\ g'',\ h'')\,(x, y, z)^2\cdot\zeta^2\\
   &\quad{} + (A,\ B,\ C,\ F,\ G,\ H)\,(x, y, z)^2\cdot 2\eta\zeta\\
   &\quad{} + (A',\ B',\ C',\ F',\ G',\ H')\,(x, y, z)^2\cdot 2\zeta\xi\\
   &\quad{} + (A'',\ B'',\ C'',\ F'',\ G'',\ H'')\,(x, y, z)^2\cdot 2\xi\eta,
\end{aligned}
\]
ou, ce qui est la m\^eme chose,
\[
\begin{aligned}
   &(a\,x^2 + b\,y^2 + c\,z^2 + 2 f\,y z + 2 g\,z x + 2 h\,x y)\,\xi^2\\
   &\quad{} + (a'\,x^2 + b'\,y^2 + c'\,z^2 + 2 f'\,y z + 2 g'\,z x + 2 h'\,x y)\,\eta^2\\
   &\quad{} + (a''\,x^2 + b''\,y^2 + c''\,z^2 + 2 f''\,y z + 2 g''\,z x + 2 h''\,x y)\,\zeta^2\\
   &\quad{} + (A\,x^2 + B\,y^2 + C\,z^2 + 2 F\,y z + 2 G\,z x + 2 H\,x y)\,2\eta\zeta\\
   &\quad{} + (A'\,x^2 + B'\,y^2 + C'\,z^2 + 2 F'\,y z + 2 G'\,z x + 2 H'\,x y)\,2\zeta\xi\\
   &\quad{} + (A''\,x^2 + B''\,y^2 + C''\,z^2 + 2 F''\,y z + 2 G''\,z x + 2 H''\,x y)\,2\xi\eta.
\end{aligned}
\]

\medskip

\noindent Soit 123 le d\'eterminant
\[
\begin{vmatrix} \partial_{\xi_1}, & \partial_{\eta_1}, & \partial_{\zeta_1}\\ \partial_{\xi_2}, & \partial_{\eta_2}, & \partial_{\zeta_2}\\ \partial_{\xi_3}, & \partial_{\eta_3}, & \partial_{\zeta_3}\end{vmatrix}
\]
et $1'2'3'$ le d\'eterminant
\[
\begin{vmatrix} \partial_{x_1}, & \partial_{y_1}, & \partial_{z_1}\\ \partial_{x_2}, & \partial_{y_2}, & \partial_{z_2}\\ \partial_{x_3}, & \partial_{y_3}, & \partial_{z_3}\end{vmatrix}
\]
et soient $U_1, U_2, U_3$ ce que devient $U$ lorsqu'on y \'ecrit $x_1, y_1, z_1, \xi_1, \eta_1, \zeta_1; x_2, y_2, z_2, \xi_2, \eta_2, \zeta_2; x_3, y_3, z_3, \xi_3, \eta_3, \zeta_3$ au lieu de $x, y, z, \xi, \eta, \zeta$; on a alors cet invariant (analogue \`a l'invariant quadratique d'une fonction binaire d'ordre pair), savoir:
\[
   \frac{1}{384}\,\overline{123}^3\cdot\overline{1'2'3'}^3\,U_1 U_2 U_3,
\]
o\`u, comme \`a l'ordinaire les valeurs $x, y, z, \xi, \eta, \zeta$ sont r\'etablies apr\`es les diff\'erentiations au lieu de $x_1, y_1, z_1, \xi_1, \eta_1, \zeta_1$, etc.

\medskip

\noindent Je repr\'esente cet invariant par
\[
\left\{\begin{array}{llllll} a, & b, & c, & f, & g, & h\\ a', & b', & c', & f', & g', & h'\\ a'', & b'', & c'', & f'', & g'', & h''\\ A, & B, & C, & F, & G, & H\\ A', & B', & C', & F', & G', & H'\\ A'', & B'', & C'', & F'', & G'', & H''\end{array}\right\}.
\]

\newpage
\setcounter{page}{351}

\noindent\textbf{271]}\hfill\textsc{sur l'invariant le plus simple d'une fonction \&c.}\hfill\textbf{351}

\medskip

\noindent on voit sans peine que l'invariant peut \^etre d\'evelopp\'e, et qu'alors il prend la forme
\[
\left\{\begin{aligned} & A,\ B,\ C,\ F,\ G,\ H\\ & a,\ b,\ c,\ f,\ g,\ h\\ & a',\ b',\ c',\ f',\ g',\ h'\\ & a'',\ b'',\ c'',\ f'',\ g'',\ h''\end{aligned}\right\}
\]
\[
\quad + 2\left\{\begin{aligned} & A',\ B',\ C',\ F',\ G',\ H'\\ & a',\ b',\ c',\ f',\ g',\ h'\\ & a,\ b,\ c,\ f,\ g,\ h\end{aligned}\right\} - \left\{\begin{aligned} & A,\ B,\ C,\ F,\ G,\ H\\ & a',\ b',\ c',\ f',\ g',\ h'\\ & a',\ b',\ c',\ f',\ g',\ h'\end{aligned}\right\}
\]
\[
\quad - \left\{\begin{aligned} & A'',\ B'',\ C'',\ F'',\ G'',\ H''\\ & a'',\ b'',\ c'',\ f'',\ g'',\ h''\\ & a,\ b,\ c,\ f,\ g,\ h\end{aligned}\right\}
\]
o\`u le premier terme
\[
\left\{\begin{aligned} & A,\ B,\ C,\ F,\ G,\ H\\ & a,\ b,\ c,\ f,\ g,\ h\\ & a',\ b',\ c',\ f',\ g',\ h'\\ & a'',\ b'',\ c'',\ f'',\ g'',\ h''\end{aligned}\right\}
\]
d\'enote la fonction
\[
\begin{aligned}
   &a b' c'' + a b'' c' + a' b'' c + a' b c'' + a'' b c' + a'' b' c\\
   &\quad{} - 2 a f' f'' - 2 a' f' f - 2 a'' f f'\\
   &\quad{} - 2 b g' g'' - 2 b' g'' g - 2 b'' g g'\\
   &\quad{} - 2 c h' h'' - 2 c' h'' h - 2 c'' h h'\\
   &\quad{} + 2 f g' h'' + 2 f g'' h' + 2 f' g'' h + 2 f' g h'' + 2 f'' g h' + 2 f'' g' h
\end{aligned}
\]
et de m\^eme pour les autres termes. L'invariant contient les termes
\[
   a b' c'' + 2 A B' G'' + 2 f g' h'' + 2 F G' H'' - 2 a f' f'' - 2 a B G + 2 a F^2 - 4 A F F'' + 4 f A F - 4 f G H,
\]

\newpage
\setcounter{page}{352}

\noindent\textbf{352}\hfill\textsc{sur l'invariant le plus simple d'une fonction \&c.}\hfill\textbf{[271}

\medskip

\noindent et on d\'eduit de l\`a son expression compl\`ete en y \'ecrivant
\[
   a b' c'' + a b'' c' + a' b'' c + a' b c'' + a'' b c' + a'' b' c
\]
au lieu de $a b' c''$;
\[
   a f' f'' + a' f'' f + a'' f f' + b g' g'' + b' g'' g + b'' g g' + c h' h'' + c' h'' h + c'' h h'
\]
au lieu de $a f' f''$, et ainsi de suite; la somme des coefficients est $6 + 2\cdot 6 + 2\cdot 6 + 4\cdot 6 + 2\cdot 9 + 2\cdot 9 + 2\cdot 9 + 4\cdot 9 + 4\cdot 9 + 4\cdot 9 = 216$, ce qui est exact, car il n'y a pas de termes qui se d\'etruisent, et dans l'expression $\tfrac{1}{384}\,\overline{123}^3\cdot\overline{1'2'3'}^3\,U_1 U_2 U_3$, le nombre des termes qui composent le produit des deux d\'eterminants carr\'es est $36\times 36 = 1296$, chacun de ces termes est multipli\'e par le facteur $8\times 8 = 64$ introduit par la diff\'erentiation, et l'on a
\[
   \frac{1}{384}\times 64\times 1296 = 216.
\]

\medskip

\noindent L'invariant qui vient d'\^etre donn\'e est li\'e avec l'invariant $T$ de M.~Aronhold (lequel se rapporte \`a une fonction ternaire cubique). C'est pour cela que je le d\'esigne par la m\^eme lettre $T$, et pour fixer sa valeur j'\'ecris
\[
   2 T = a b' c'' + \text{etc.}
\]

\medskip

\noindent Cela pos\'e, je forme les deux combinants $C_6$ et $C_{12}$ de la mani\`ere suivante.

\medskip

\noindent\textit{Combinant du sixi\`eme ordre $C_6$.}\ \ Soient $U, V, W$ les trois fonctions quadratiques ternaires; consid\'erez la fonction syzyg\'etique $\lambda U + \mu V + \nu W$, o\`u $\lambda, \mu, \nu$ sont des quantit\'es arbitraires; formez avec les variables $\xi, \eta, \zeta$ le r\'eciprocant\footnote{Le r\'eciprocant de $(a, b, c, f, g, h)\,(x, y, z)^2$ est
\[
   \left|\begin{array}{cccc} \xi, & \eta, & \zeta\\ a, & h, & g, & \xi\\ h, & b, & f, & \eta\\ g, & f, & c, & \zeta\end{array}\right| = -\,(\mathfrak{A},\mathfrak{B},\mathfrak{C},\mathfrak{F},\mathfrak{G},\mathfrak{H})(\xi, \eta, \zeta)^2.
\]
} (multipli\'e par 2) de cette fonction, le r\'esultat sera de l'ordre 2 par rapport aux coefficients de $U, V, W$, de l'ordre 2 par rapport \`a $\lambda, \mu, \nu$, et de l'ordre 2 par rapport \`a $\xi, \eta, \zeta$; c.\ \`a d.\ ce r\'esultat sera une fonction quadratique bi-ternaire de $\lambda, \mu, \nu$ et $\xi, \eta, \zeta$ dont l'invariant $2 T$ (lequel sera par cons\'equent de l'ordre 6 par rapport aux coefficients de $U, V, W$) est le combinant $C_6$ qu'il s'agissait de trouver.

\medskip

\noindent\textit{Combinant du douzi\`eme ordre $C_{12}$.}\ \ Consid\'erez comme auparavant la fonction syzyg\'etique $\lambda U + \mu V + \nu W$, formez le discriminant (multipli\'e par 6) de cette fonction; le r\'esultat sera de l'ordre 3 par rapport aux coefficients de $U, V, W$, et de l'ordre 3 par rapport \`a $\lambda, \mu, \nu$; c.\ \`a d.\ il sera une fonction cubique ternaire de $\lambda, \mu, \nu$, dont

\newpage
\setcounter{page}{353}

\noindent\textbf{271]}\hfill\textsc{sur l'invariant le plus simple d'une fonction \&c.}\hfill\textbf{353}

\medskip

\noindent l'invariant $S$ de M.~Aronhold\footnote{Pour la fonction cubique ternaire $(a, b, c, i, j, k, i_1, j_1, k_1, l)\,(x, y, z)^3$, l'invariant $S$ est
\[
\begin{aligned}
   &-\,2 l^2\,(k j_1 + k_1 j + i_1 j_1)\\
   &\quad{} + 3 l\,(i j k + i_1 j_1 k_1)\\
   &\quad{} + l\,(a j_1 j + b k_1 i_1 + c i_1 j_1)\\
   &\quad{} - l\,a b c\\
   &\quad{} - (a i j_1 k_1 + b j k_1 i_1 + c k i_1 j_1)\\
   &\quad{} + (j^2 k_1^2 + k^2 i_1^2 + i^2 j_1^2)\\
   &\quad{} + (b c j_1 k_1 + c a k_1 i_1 + a b i_1 j_1)\\
   &\quad{} - (a j_1^2 j + b k_1^2 k + c i_1^2 i + a i^2 k_1 + b j^2 i_1 + c k^2 j_1).
\end{aligned}
\]
} (lequel sera de l'ordre 12 par rapport aux coefficients de $U, V, W$) est le combinant $C_{12}$ qu'il s'agissait de trouver.

\medskip

\noindent On peut exprimer ces r\'esultats d'une mani\`ere abr\'eg\'ee en disant que $C_6$ est l'invariant $2 T$ de la fonction bi-ternaire $2 R$ et que $C_{12}$ est l'invariant $S$ de la fonction cubique ternaire $6 D$, en d\'esignant par $R$ et $D$ le r\'eciprocant et le discriminant de la fonction $\lambda U + \mu V + \nu W$.

\medskip

\noindent Pour d\'emontrer l'\'equation
\[
   12 R = 16 C_{12} - C_6{}^2
\]
il suffit de faire voir que cette \'equation est v\'erifi\'ee lorsqu'au lieu des trois fonctions donn\'ees $U, V, W$ on substitue trois fonctions de la forme $l U + m V + n W$, les coefficients $l, m, n$ ayant \'et\'e choisis de mani\`ere \`a simplifier autant que possible les expressions des trois fonctions. Il est facile de voir qu'il est permis de prendre pour les trois fonctions les formes dont s'est servi M.~Sylvester dans son m\'emoire, savoir
\[
\begin{aligned}
   &p\,(\xi^2 - \zeta^2),\\
   &q\,(\eta^2 - \zeta^2),\\
   &y^2 + 2 f y z + 2 g z x + 2 h x y.
\end{aligned}
\]
Mais il suit des recherches de M.~Hesse, et M.~Sylvester aussi a depuis reconnu, qu'on peut prendre pour les trois fonctions les formes encore plus simples
\[
\begin{aligned}
   &x^2 + 2 l y z,\\
   &y^2 + 2 l z x,\\
   &z^2 + 2 l x y,
\end{aligned}
\]
qui se r\'eduisent aux fonctions d\'eriv\'ees de la fonction cubique $x^3 + y^3 + z^3 + 6 l x y z$. Or, en prenant ces valeurs de $U, V, W$, on obtient
\[
   \lambda U + \mu V + \nu W = (\lambda, \mu, \nu, \lambda l, \mu l, \nu l)\,(x, y, z)^2.
\]

\hfill C.~IV.\hfill 45

\selectlanguage{english}

\end{document}
